% Build from this folder with tectonic final_report.tex or two pdfLaTeX passes.
\documentclass[11pt,a4paper]{article}
\usepackage[T1]{fontenc}
\usepackage{lmodern}
\usepackage{upquote}
\usepackage{microtype}
\usepackage[top=3cm,bottom=3cm,left=2cm,right=2cm]{geometry}
\usepackage{amsmath,amssymb}
\usepackage{array,booktabs,tabularx}
\usepackage{graphicx}
\usepackage{fancyhdr}
\usepackage[table]{xcolor}
\usepackage{caption}
\captionsetup{hypcap=false,font=small,labelfont=bf,justification=raggedright,singlelinecheck=false,skip=4pt}
\captionsetup[figure]{position=bottom}
\captionsetup[table]{position=bottom}
\newcommand{\reportfigure}[5][]{%
  \begin{center}\begin{minipage}{\linewidth}\centering
  \includegraphics[width=#2]{#3}\par
  \captionof{figure}[#1]{#4}\label{#5}%
  \end{minipage}\end{center}%
}
\usepackage[hidelinks]{hyperref}
\definecolor{LinkInk}{HTML}{154C62}\newcommand{\reftitle}[2]{\href{#1}{\textcolor{LinkInk}{\textquotedblleft#2\textquotedblright}}}

\newcommand{\UniversityLineOne}{VIETNAM NATIONAL UNIVERSITY HO CHI MINH CITY}
\newcommand{\UniversityLineTwo}{HO CHI MINH CITY UNIVERSITY OF TECHNOLOGY}
\newcommand{\FacultyName}{FACULTY OF COMPUTER SCIENCE AND ENGINEERING}
\newcommand{\ReportCourse}{Specialized Project (CO4029)}
\newcommand{\ReportTitle}{Budget-Aware Densification for Efficient 3D Gaussian Splatting: Surface Reconstruction and Editable Mesh Evaluation}
\newcommand{\LecturerName}{Nguyen Tuan Khoi}
\newcommand{\StudentName}{Nguyen Vu Long}
\newcommand{\LogoFile}{figures/hcmut-logo.png}

\definecolor{ReportBlue}{HTML}{123D70}
\definecolor{ReportLight}{HTML}{EDF3F8}
\definecolor{ReportGray}{HTML}{526170}
\hypersetup{pdftitle={Budget-Aware Densification for Efficient 3D Gaussian Splatting},
  pdfauthor={Nguyen Vu Long},pdfsubject={Specialized Project (CO4029)}}
\renewcommand{\contentsname}{TABLE OF CONTENTS}
\setcounter{tocdepth}{2}
\setlength{\headheight}{40pt}
\setlength{\parindent}{0pt}
\setlength{\parskip}{5pt}
\setlength{\emergencystretch}{1em}
\renewcommand{\arraystretch}{1.13}
\newcolumntype{Y}{>{\raggedright\arraybackslash}X}
\makeatletter
% Read generated alignment rows directly, without the file hooks of LaTeX's \input.
\newcommand{\inputrows}[1]{\@@input #1 }
\makeatother
\raggedbottom

\pagestyle{fancy}
\fancyhf{}
\fancyhead[L]{\begin{tabular}{@{}c@{\hspace{7pt}}l@{}}
  \raisebox{-0.40\height}{\includegraphics[width=8mm]{\LogoFile}} &
  \begin{tabular}{@{}l@{}}\footnotesize\bfseries\ttfamily\UniversityLineTwo\\
  \footnotesize\bfseries\ttfamily\FacultyName\end{tabular}
\end{tabular}}
\fancyfoot[L]{\scriptsize\ttfamily CO4029 / HK261}
\fancyfoot[R]{\scriptsize\ttfamily Page \thepage/\pageref*{report-end}}
\renewcommand{\headrulewidth}{0.3pt}
\renewcommand{\footrulewidth}{0.3pt}

\setlength{\topsep}{3pt}\begin{document}
\hypersetup{pageanchor=false}
\begin{titlepage}
\centering
{\bfseries\UniversityLineOne\\[1mm]
\UniversityLineTwo\\[1mm]
\FacultyName\par}
\vspace{1cm}
\includegraphics[width=3cm]{\LogoFile}\par
\vspace{1cm}

\begin{minipage}{0.88\textwidth}
\raggedright
{\Large\bfseries\ReportCourse\par}
{\large Final Report\par}
\vspace{0.55cm}
\hrule
\vspace{0.45cm}
{\Large\bfseries\ReportTitle\par}
\vspace{0.45cm}
\hrule
\vspace{0.55cm}
\begin{tabular}{@{}ll@{}}
\textbf{Lecturer} & \LecturerName \\
\textbf{Student} & \StudentName{} (2352696) \\
\textbf{Email} & \href{mailto:long.nguyenty@hcmut.edu.vn}{long.nguyenty@hcmut.edu.vn} \\
\end{tabular}
\end{minipage}

\vspace{2cm}
{\small\bfseries Ho Chi Minh City, 20 November 2026\par}
\end{titlepage}
\hypersetup{pageanchor=true}


\section*{Abstract}
\addcontentsline{toc}{section}{Abstract}
This project tests whether a 3D Gaussian Splatting model can use fewer primitives while preserving image quality and yielding an editable mesh. It compares the default gsplat strategy (B00) with A05-Abs, a policy that selects Gaussians for splitting using absolute image-space gradients and a scheduled count ceiling. At 7,000 matched updates, A05-Abs used 45.57--54.56\% fewer Gaussians across four public indoor development scenes and raised mean held-out PSNR by 0.210--0.739 dB. The Bonsai mean depended strongly on one view. Across three Room seeds, the candidate gained an average of 0.674 dB and used 45.6\% fewer primitives. The 700,000 ceiling admitted every eligible parent at the selected Room threshold, while the same ceiling bound all refinement events on the outdoor Garden scene, where the candidate gained 0.365 dB with 75.6\% fewer Gaussians. On a local RealSense D455 segment, the candidate used 62.94\% fewer Gaussians and its mesh's median absolute sensor-depth discrepancy rose from 8.38 to 9.47 cm. All ten raw-mesh edits and the cleaned Room tabletop lift persisted after reopening. Meshes extracted by fusing rendered depth retained open boundaries. The results support a rendering-quality and size advantage on the tested public scenes and demonstrate local editability. Independent physical surface accuracy remains unresolved.\par
\textbf{Keywords} 3D Gaussian Splatting, adaptive densification, RGB-D reconstruction, volumetric fusion, mesh editing\par
\clearpage
\begingroup\fontsize{9.5}{11}\selectfont\setlength{\parskip}{0pt}\tableofcontents\endgroup
\clearpage
\phantomsection\addcontentsline{toc}{section}{List of Figures}
\listoffigures
\clearpage
\phantomsection\addcontentsline{toc}{section}{List of Tables}
\listoftables
\clearpage
\section{Introduction}
\subsection{Motivation and problem statement}
Photographic reconstruction serves two related tasks. In view synthesis, the system predicts what a scene looks like from another camera. In geometric reconstruction, it estimates where the visible surfaces lie and how they connect. A virtual walkthrough mainly benefits from realistic views and interactive rendering. An editing workflow also needs a polygon asset whose vertices, boundaries and connected regions can be selected and changed. A system that serves both tasks must account for the different evidence required to evaluate them.

Three-dimensional Gaussian Splatting (3DGS) is useful for view synthesis because its explicit primitives are easy to render and can be allocated where the fitted representation needs them \cite{ref4}. The optimization changes both the attributes and the number of these primitives. This creates a resource-allocation problem. More splats require more parameter storage and optimizer state. They also affect the number of projected footprints and their overlap in the renderer. Adding a primitive only pays off when it improves the representation enough to justify these costs. The magnitude of the benefit depends on the scene, the training stage and the allocation criterion.

Converting the resulting representation to polygons adds a second problem. A Gaussian radiance field assigns color and opacity over a volume. The observed RGB image constrains their composite. A mesh requires an explicit choice about the location of the surface and the connectivity of neighboring samples. In the public Room photographs, a coffee table has a broad top, clutter and narrow supports. The renderer can reproduce the tabletop color while the extracted surface joins clutter to the top. The editor then inherits this join. Figure~\ref{figure1} traces one tabletop example through the pipeline.

\reportfigure[From a Room photograph to a reopened tabletop edit]{0.99\linewidth}{figures/teaser.jpg}
  {Representative public Room result. Left to right, held-out photograph, A05-Abs Gaussian RGB render, extracted mesh and reopened tabletop edit. The first three panels use the same tabletop crop. The ring marks the raw-mesh edit target. The last panel shows a separate edit on the cleaned surface.}{figure1}

This work studies allocation first and tests what the resulting representation supports downstream. A05-Abs combines a count schedule, absolute screen-space positional gradients, image-based ranking and Long-Axis Split inside the existing gsplat training pipeline. The software exports the optimized models and converts their rendered depths into raw meshes. A controlled photo-only benchmark tests representation efficiency. A bounded RGB-D capture adds a same-device depth comparison. Both conditions use a reproducible mesh-editing task.

\subsection{Research questions and success criteria}
The primary question concerns the trade-off between rendering quality and representation size. At the same update count, can the method achieve equal or better mean held-out RGB quality with fewer Gaussian primitives than the local default trainer? The primary evidence is the paired PSNR and final count per public scene. SSIM and LPIPS indicate whether the direction depends on the image metric. Checkpoint size, tensor allocation and render-call throughput provide separate resource measurements. A matched elapsed-time study compares the two models at common time limits. The Room sweep measures the realized operating points when the gradient threshold or available capacity changes.

The second question concerns the mechanism. How sensitive are those outcomes to the gradient definition, the numerical threshold, the error and edge ranking terms, and the split operator? The combined comparison estimates the effect of changing the whole strategy. The component experiments perturb a smaller part of the strategy with the initializer, camera order, objective and update count fixed. For the earlier A05-Abs study against A04, the practical criterion was a PSNR increase of at least 0.05 dB with SSIM degradation and LPIPS increase each at most 0.001. The new Room screen reports measured effect sizes and count changes for its individual contrasts. The same scene was used during development, so these contrasts diagnose the implementation on that scene.

The third question is whether the compact representation preserves the properties useful for meshing. A successful conversion produces finite vertices and faces, preserves scene coordinates and exports both PLY and OBJ. Component counts, boundary edges and inspected crops characterize fragmentation and detail. On the matte RGB-D segment, camera-axis agreement is measured using reserved depth messages and common valid pixels. Lower absolute discrepancy at a comparable coverage would support improved consistency on that segment. Absolute dimensional accuracy requires an external reference and is outside the evidence collected here.

Table~\ref{table1} summarizes the four comparisons. The fourth question concerns local editability. The operational criterion is a connected patch with nonzero and measurable displacement after a separate process reopens the saved Blender project. Mesh topology, watertightness and accuracy have their own evidence requirements. A file that passes the edit can still contain open boundaries, disconnected fragments or merged thin parts. The experiment records these outcomes separately.

\begin{center}\begin{minipage}{\linewidth}
\small\begin{tabularx}{\linewidth}{@{}p{.24\linewidth}YY@{}}
\toprule
\textbf{Question} & \textbf{Controlled comparison} & \textbf{Measured outcome}\\
\midrule
Representation efficiency & Matched fitting order, targets and optimization horizon & Gaussians, RGB metrics, storage, tensor allocation and render time\\
Allocation mechanism & Gradient-by-threshold factorial, ranking ablation, cap sweep and earlier split factorial & Paired changes in RGB metrics and final count\\
Surface reconstruction & Shared fitting cameras and TSDF settings. Separate grid-resolution sensitivity & Coverage, topology, inspected details and paired sensor-depth discrepancy\\
Local design editing & Connected patch selected and measured after reopening & Persisted displacement, evaluated mesh and file hash\\
\bottomrule
\end{tabularx}\par
\captionof{table}[Research questions, comparisons and outcome measures]{Research questions and the evidence used to answer them.}\label{table1}
\end{minipage}\end{center}

\subsection{Scope and requirements}
The main benchmark contains four static or mostly static public indoor scenes from Mip-NeRF 360 \cite{ref2}. Each uses the supplied structure-from-motion model, factor-two photographs and the same split rule across arms. Public Room is the main case for controlled extensions and the complete tabletop example. A separate outdoor Garden test applies the frozen policy to factor-four images from the same dataset family. The museum comparison uses a short RGB-D segment from a personal moving-camera recording. Its three nearby check frames and shared sensor calibration constrain the interpretation of depth disagreement.

The system is designed for a single 6 GiB desktop GPU and approximately 16 GiB of system memory. The training configuration keeps a batch of one camera, uses packed rendering and uses a prepared initializer. A cap on live primitives sets one explicit resource constraint. Working tensors, tile intersections, cached images and exported artifacts have separate memory and storage costs. The extractor uses a bounded dense volume and reduces the spatial resolution of its fitting-camera renders. These choices support repetition on the project workstation.

The public route starts with prepared poses. It measures representation and meshing given that camera model. A separate video route diagnoses camera recovery before training, and the museum route tests local camera fitting with stronger depth provenance. The thesis uses those acquisition trials to state when a matched reconstruction is possible. It treats full-route museum reconstruction as deferred because the local validation cannot establish a globally reliable model.

The project preserves raw source and intermediate outputs alongside derived assets. Rough geometry is useful for diagnosing camera and representation failures. A visually cleaner derivative receives its own record so it can be compared with the raw mesh and reprocessed from the same input. Evaluation source identities and hashes are recorded before pairwise comparisons. This allows an examiner to identify changes introduced by training, extraction and editing separately.

\subsection{Contributions and report organization}
The first contribution is a documented allocation policy in the native gsplat loop. Its project-specific choices are the event-indexed ceiling, candidate ranking and support discount, integration of the existing absolute gradients, and use of Long-Axis Split for all selected parents. Splatting, absolute-gradient accumulation, Long-Axis Split and volumetric fusion draw on prior work. Section~\ref{sec:related} specifies the relationship to those methods.

The second contribution is a matched evaluation from prepared cameras through trained models and raw meshes. The public comparisons pair initial weights, camera order and validation RGB targets. Component and resource studies on Room examine which results depend on the allocation design or operating point. The matte comparison separates the agreement of Gaussian expected depth from the agreement of mesh first-hit depth. Fixed-checkpoint extractions probe sensitivity to TSDF resolution.

The third contribution is an executable editing workflow backed by independent Blender processes. It records raw and cleaned geometry, a selected edit region and displacement in the reopened file. This workflow addresses a practical handoff from a trained radiance field to an editable polygon asset.

Section 2 develops the camera, radiance field and fusion models. Section 3 gives the training and refinement algorithms and the data contracts between stages. Section 4 presents the paired results, component experiments and the validity limits of each measurement. Section 5 answers the research questions against that evidence. The appendices provide reproduction commands and source records.

\clearpage
\section{Theoretical Background and Related Work}
\subsection{Camera projection and coordinate frames}
Let $X_w$ be a point in world coordinates. A calibrated camera stores the transformation from world to camera coordinates as $X_c=R_{cw}X_w+t_{cw}$. The translation in this expression is the position of the world origin in camera coordinates. The camera center in the world is $-R_{cw}^{\top}t_{cw}$. For the undistorted pinhole case, with $X_c=(x,y,z)^\top$ and $z>0$, the projective relation is
\begin{equation}
  z\begin{bmatrix}u\\v\\1\end{bmatrix}
  = KX_c,\qquad
  K=\begin{bmatrix}f_x&0&c_x\\0&f_y&c_y\\0&0&1\end{bmatrix}.
  \label{eq:project}
\end{equation}
Here $u=f_x x/z+c_x$ and $v=f_y y/z+c_y$. The intrinsics set the mapping between image and camera coordinates. The extrinsics set the mapping between cameras and the world. Distorted cameras require a calibrated distortion model or prior undistortion. When image dimensions change, the corresponding rows of $K$ scale with the width and height ratios. Mixing resolutions without that adjustment changes the ray direction for the same imaged feature.

Backprojection at a known camera-axis depth $z$ gives $X_c=zK^{-1}(u,v,1)^\top$. The last component of the unnormalized ray direction is one. The ray parameter is then camera $z$ even away from the principal point. A unit direction instead parameterizes distance along the ray, which is larger than $z$ away from the center. The evaluation uses camera-axis depth throughout. For relative transforms between two cameras, the source sample is backprojected, moved through the relative rigid transform, then reprojected with the target camera intrinsics. This gives the geometric basis of sensor-to-color alignment and track-based checks.

Under a camera similarity transform a normalized position has the form $X'=aQX+b$ with $a>0$. Applying this transformation to both the point cloud and camera centers and rotating their coordinate axes preserves the perspective projections. Camera-axis distances scale by $a$. Public SfM is defined up to a similarity. Normalized units therefore support within-model comparisons and reproducible edit offsets; converting an offset to centimeters requires a calibrated scale. In the matte experiment, the metric depth values help estimate the map scale before the loader normalizes the map. The evaluation divides rendered axis depths by the recorded normalization scale to restore sensor units.

\subsection{Structure from Motion, bundle adjustment and depth}
Structure from Motion estimates camera poses and a sparse 3D structure from overlapping images \cite{ref1}. A feature track links measurements predicted to arise from the same world point. Given an initial relative pose, triangulation finds a point that approximately agrees with those viewing rays. Given initial camera poses and point positions, bundle adjustment refines both by minimizing reprojection error over the observed tracks. With observed feature $u_{ij}$ for point $j$ in camera $i$, a common formulation is
\begin{equation}
  \min_{\{R_i,t_i,X_j\}}\sum_{(i,j)\in\mathcal O}
        \rho\bigl(\|\pi(K_i(R_iX_j+t_i))-u_{ij}\|_2^2\bigr),
  \label{eq:ba}
\end{equation}
where $\pi$ divides the first two homogeneous coordinates by the third. A robust loss $\rho$ limits the influence of large residuals, and an explicit gauge fixes arbitrary global degrees of freedom. The normal equations are sparse because an observation couples one camera and one track. Nearly parallel rays and short baselines can leave poor depth conditioning even when reprojection error is small. Repeated texture can create mutually consistent appearance matches at the wrong locations. Moving people, specular regions and occlusion can violate the assumption of a fixed 3D point.

An aligned RGB-D observation adds a noisy measurement $z_{ij}^{\rm sensor}$ of camera-axis depth. A metric residual compares it to $\mathbf e_3^\top(R_iX_j+t_i)$, scaled by an uncertainty estimate for that observation. A long-range or depth-discontinuity sample can be less informative than a nearby pixel on a stable matte surface. Confidence weights, robust loss and held-out correspondences control different failure modes. The first two govern the fit. The held-out correspondences test consistency of the fitted map on measurements excluded from that fit. The project also reserves depth messages at validation cameras when measuring model depth. Their RGB can localize the camera against the frozen map, so the measurements share the acquisition device and a reference coordinate system.

Aligned depth uses calibration and synchronization in addition to the measurement itself. Given a raw point in the depth camera, the relative rotation and translation map it to the color camera. Projection to color pixels may map several samples to one pixel or leave holes. Calibration errors, scene motion between exposures and mixed optical returns can bias the aligned depth. Fusing more registered observations helps only when their errors and coverage support the same surface. The local matte tests evaluate this restricted case as described in Section~\ref{sec:acquisition}.

\subsection{From volume rendering to explicit Gaussian primitives}
NeRF parameterizes density $\sigma$ and directional color $c$ with a continuous field \cite{ref3}. For ray $r(t)$ and near/far range $[t_n,t_f]$, the foreground color contribution is
\begin{equation}
  C(r)=\int_{t_n}^{t_f}\!T(t)\,\sigma(r(t))\,c(r(t),\mathbf d)\,dt,\qquad
  T(t)=\exp\!\left(-\int_{t_n}^{t}\!\sigma(r(s))\,ds\right).
\end{equation}
Discretizing a segment of length $\Delta_i$ with constant $\sigma_i$ gives opacity $1-\exp(-\sigma_i\Delta_i)$. Multiplying by the transmittance through the preceding segments yields the segment weight. Both NeRF and 3DGS use this compositing logic. Their representations differ. NeRF repeatedly queries a field along rays. 3DGS stores explicit primitives and rasterizes their projected footprints \cite{ref4}. The discrete representation allows directly splitting, cloning, pruning and counting primitives.

For Gaussian $i$, let $\mu_i\in\mathbb R^3$ be its center, $R_i$ the rotation from its normalized quaternion, and $S_i=\mathrm{diag}(s_{i1},s_{i2},s_{i3})$ its positive scales. A point at local coordinate $y$ is mapped by $X=\mu_i+R_iS_i y$. If the covariance of $y$ is $I$, the transformed covariance is
\begin{equation}
 \Sigma_i=R_iS_iS_i^\top R_i^\top,\qquad
 g_i(X)=\exp\!\left[-\tfrac12(X-\mu_i)^\top\Sigma_i^{-1}(X-\mu_i)\right].
\end{equation}
This convention describes a unit-height Gaussian shape. In training, log-scales are exponentiated and opacity logits are passed through a sigmoid. These parameterizations keep scales positive and activated opacity between zero and one. The renderer evaluates spherical harmonic coefficients at the camera viewing direction to permit view-dependent colors. There are 16 real basis values per color channel at degree three. The feature tensor stores its constant coefficient separately from the rest.

For an intuition for parameter storage, the standard RGB case needs three means, four quaternion entries, three log-scales, one opacity logit and 48 color coefficients per splat. With 32-bit floating-point parameters this is 236 bytes per splat before checkpoint metadata. Two Adam moment tensors can double the parameter-sized part of the optimizer state. Forward and backward buffers add further memory and depend on resolution, overlap and the rendering implementation. Count thus measures one part of the cost, while peak allocated memory and render time measure properties of an actual run.

\subsection{Projection, alpha compositing and expected depth}
Perspective projection is nonlinear. Let $(x,y,z)$ now denote a Gaussian center in camera coordinates. Linearization of Equation~\ref{eq:project} around the center gives the Jacobian
\begin{equation}
 J=\begin{bmatrix}
 f_x/z &0&-f_xx/z^2\\
 0&f_y/z&-f_yy/z^2
 \end{bmatrix}.
\end{equation}
To first order, perturbations are transformed by $JR_{cw}$, so the footprint covariance is $JR_{cw}\Sigma_i R_{cw}^{\top}J^\top$, up to the rasterizer's footprint treatment. The off-diagonal terms determine the orientation of the ellipse. At a pixel, the model multiplies the Gaussian's learned opacity by its evaluated footprint. The CUDA rasterizer bins screen-space footprints into tiles and composites covered pixels in depth order. It evaluates only projected entries that can contribute, and the packed interface records their Gaussian IDs \cite{ref5}.

Let $\alpha_i(p)$ denote the effective alpha and $c_i$ the evaluated color from primitive $i$ in one camera. For front-to-back contributors,
\begin{align}
 T_i(p)&=\prod_{j<i}(1-\alpha_j(p)),\qquad w_i(p)=T_i(p)\alpha_i(p),\\
 A(p)&=\sum_i w_i(p)=1-\prod_i(1-\alpha_i(p)),\\
 C(p)&=\sum_i w_i(p)c_i+(1-A(p))C_{\rm bg},\\
 D_{\rm ED}(p)&=\frac{\sum_i w_i(p)z_i}{A(p)}\qquad (A(p)>0).
 \label{eq:expected-depth}
\end{align}
The accumulated alpha equality follows from $T_i-T_{i+1}=T_i\alpha_i$. Expected depth normalizes by accumulated alpha and uses the projected center depths $z_i$. Even at high alpha, a ray can average several layers (Figure~\ref{figure2}). If two weights are 0.60 and 0.30 at depths 1 and 2, the expected depth is $1.33$ despite the nearer contributing layer at 1. Other combinations of opacity and color may produce similar RGB under the available views. The measured RGB objective leaves this ambiguity at mixed pixels.

\reportfigure[Expected depth along a ray with two contributing Gaussians]{0.84\linewidth}{figures/depth_illustration.png}
  {Projection and depth ambiguity in Gaussian rendering. The two ellipsoids both contribute to one ray. Their normalized expected depth falls between them.}{figure2}

Mesh evaluation uses the first positive triangle intersection along the camera ray. This value is a visibility decision made on the already extracted polygon surface. It is different from the opacity-weighted average on the Gaussian model. A first hit can move toward the camera if fusion bridges a gap with a foreground surface. It can move away from the camera if an expected front surface is missing. Evaluating both quantities against the same sensor source can distinguish disagreement present before extraction from changes associated with the downstream process. Shared support is essential because missing hits alter which surfaces remain in the sample.

\subsection{Densification and gradient cancellation}
Training alternates continuous optimization of Gaussian parameters with changes to the primitive population. In 3DGS, a high average positional gradient indicates that moving a projected mean may reduce the image objective. The default strategy uses this signal alongside a size threshold to choose cloning or splitting \cite{ref4,ref5}. It also prunes least opaque primitives and limits excessively large ones under its pruning schedule. This feedback creates capacity around areas with fitting pressure, while the final count emerges from the optimization.

For projected mean $m_{iv}$ of primitive $i$ in view $v$, let $b_{ivp}$ denote the contribution of pixel $p$ to the gradient of the objective with respect to that mean. The signed renderer returns $\sum_p b_{ivp}$. If two pixels yield $(g,0)$ and $(-g,0)$, the net is zero even though both pixels exert fitting pressure. AbsGS aggregates component-wise magnitudes before taking a vector norm \cite{ref6}. In the example it yields $(2|g|,0)$. More generally the triangle inequality and monotonicity of the norm imply
\begin{equation}
 \left\|\sum_p b_{ivp}\right\|_2
  \ \leq\ \left\| \sum_p |b_{ivp}| \right\|_2 ,
\end{equation}
where the absolute value is applied component by component. This argument holds for identical pixels and Gaussian parameters in one backward pass. Changing the criterion also changes the subsequent population and optimization. A threshold calibrated for signed gradients need not select the same quantity or mix of candidates under absolute accumulation.

Gradient magnitude quantifies a local derivative, while the residual quantifies how far the composite image lies from its target. Several Gaussians can share responsibility for the same residual. An edge map highlights contrast in the target image, including texture, shadows, occlusion and silhouettes. It does not identify which edges represent a physical boundary. The policy combines these cues as a ranking after applying a gradient gate. Section~\ref{sec:allocation} specifies the normalization and allocation steps. Taming 3DGS and Improving Densification provide relevant examples of budget scheduling and contribution-aware densification \cite{ref12,ref7}.

\subsection{Truncated signed distance fusion}
Given a depth image $D_v$ and a voxel center $X$, let $z_v(X)$ be its axial depth and let $p_v(X)$ be its projection in view $v$. A projective signed distance is proportional to $D_v(p_v(X))-z_v(X)$. A voxel in front of the observed surface has positive distance under this convention. Restricting updates near and in front of the observation avoids treating arbitrarily far occluded space as measured geometry. The scalar field averages observations with limited influence beyond a truncation distance $\tau$. A simplified normalized update is
\begin{align}
 F_v(X)&=\mathrm{clip}\bigl((D_v(p_v(X))-z_v(X))/\tau,-1,1\bigr),\\
 F_{\rm fused}(X)&=\frac{\sum_{v\in\mathcal V(X)}w_v(X)F_v(X)}
                         {\sum_{v\in\mathcal V(X)}w_v(X)} .
\end{align}
The contributing set $\mathcal V(X)$ excludes unsupported and occluded samples according to the integrator. The observed sign change approximates the surface by a zero crossing. Within a cell edge joining opposite signs $f_0$ and $f_1$, the linear estimate lies at a fraction $f_0/(f_0-f_1)$ of the edge. Marching Cubes generates triangles from the sign pattern at each cell's corners and interpolates their vertices along the intersected edges \cite{ref14}. This construction explains why voxel size and valid-view coverage affect connectivity, even when the input Gaussians are fixed \cite{ref9}. Open3D integrates the observations and extracts the triangles. An alternative, Poisson surface reconstruction, estimates an implicit indicator function from consistently oriented surface samples \cite{ref15}. The Gaussian centers in this project carry appearance and opacity, and their covariance axes supply no verified surface normals. Fusing camera depths gives both arms the same observable input to the extractor and allows us to measure how grid spacing affects the result.

An observation can already blur a thin feature due to depth averaging or low-resolution rendering. Fusion can broaden or suppress the feature further if the truncation band covers two nearby surfaces or if a gap spans too few voxels. Conversely, a finer grid can expose holes or tiny fragments that a coarse grid bridges. A denser mesh is therefore an ambiguous outcome. The fixed-checkpoint sweep in Section~\ref{sec:grid-results} changes grid spacing while holding the fusion cameras and truncation distance fixed. It estimates sensitivity to discretization under the existing depth sampling, crop and visibility policy.

\subsection{Related methods and the research gap}\label{sec:related}
Table~\ref{tab:related} identifies which stages reuse prior methods and which choices are evaluated here. SfM supplies consistency in image projection and sparse structure. NeRF and 3DGS supply view synthesis. AbsGS modifies the statistic used for densification. Budget-aware methods change the distribution of new primitives under resource constraints. Surface-oriented methods constrain the representation or renderer before extracting geometry.

\begin{center}\begin{minipage}{\linewidth}
\small\begin{tabularx}{\linewidth}{@{}>{\raggedright\arraybackslash}p{.20\linewidth}YY@{}}
\toprule
\textbf{Method} & \textbf{Design objective} & \textbf{Relationship to A05-Abs}\\
\midrule
COLMAP \cite{ref1} & Incremental calibration and sparse reconstruction from image correspondences & Supplies public camera models and the feature/track framework for local RGB-D pose fitting\\
NeRF \cite{ref3} & Radiance and density queried along rays & Provides a reference model for compositing and depth ambiguity\\
3DGS and gsplat \cite{ref4,ref5} & Explicit Gaussian primitives and efficient differentiable rasterization & Shared representation, rendering and optimization infrastructure. The default strategy is B00\\
AbsGS \cite{ref6} & Absolute per-pixel view-space gradients to reduce cancellation & Uses gsplat's absolute-gradient path as one input to capped candidate selection\\
Taming 3DGS \cite{ref12} & Budgeted growth using contribution cues, plus renderer optimizations & Related resource objective. The project uses a ceiling that can be underfilled\\
Improving Densification \cite{ref7} & Edge-aware allocation and Long-Axis Split, with additional pruning and update changes & Adapts its split geometry and uses a local deterministic top-$k$ ranking with a residual term\\
SuGaR and 2DGS \cite{ref8,ref13} & Surface-aligned primitives or planar splats designed for geometry & Alternatives when geometric constraints are imposed during training. This study compares matched 3DGS arms\\
Volumetric fusion \cite{ref9} & Combine range observations into an implicit surface & Fuses expected depths rendered from the learned scene after training\\
\bottomrule
\end{tabularx}\par
\captionof{table}[Prior methods and their relationship to the project]{Related methods and their roles in the project.}\label{tab:related}
\end{minipage}\end{center}

Published methods already address gradient collision, footprint placement and Gaussian budgets. The research gap addressed here is the behavior of a specific integration under matched scene and compute conditions, including the consequences for exported polygons and editing. The study retains the same renderer in both arms and measures the output at multiple stages. It can then attribute some downstream differences to the learned model given the chosen extractor, and measure extraction sensitivity given fixed learned models. It does not compare the proposed allocation method head-to-head with a separate surface-oriented renderer.

\clearpage
\section{Proposed Methodology and System Design}
\subsection{Problem formulation and architecture}
Given calibrated fitting cameras $\mathcal T$, their RGB images and an initialized point set, the training stage produces a set of optimized Gaussians $\mathcal G$. For the matte route it also receives weighted depth samples at fitting views. The renderer maps $(\mathcal G,K,T_{cw})$ to RGB, alpha and optionally expected depth. The allocation strategy controls the live size of $\mathcal G$ during refinement. The extraction stage consumes the final model and fitting cameras, and the editing stage consumes the extracted mesh. Figure~\ref{figure3} traces these stages. Table~\ref{tab:interfaces} identifies the outputs and checks at each handoff.

\reportfigure[Acquisition, training, mesh extraction and editing workflow]{0.99\linewidth}{figures/pipeline.png}
  {Acquisition, Gaussian training, volumetric extraction and polygon editing. The public route reads prepared cameras. The RGB-D route estimates a local map and reserves depth messages for the depth comparison.}{figure3}

The public experiment uses an RGB objective and a common extractor. Comparing the two learned models at identical TSDF settings controls the extraction procedure. Changing grid resolution with either model fixed probes a different source of variability. The editor operates on triangle connectivity and never feeds its modifications back to the learned model. The raw export is therefore the appropriate interface for comparing two training strategies. A cleaned mesh serves the specific downstream design task.

\begin{center}\begin{minipage}{\linewidth}
\small\begin{tabularx}{\linewidth}{@{}p{.16\linewidth}YY@{}}
\toprule
\textbf{Stage} & \textbf{Input and output} & \textbf{Check}\\
\midrule
Camera setup & Photographs and COLMAP model, or calibrated color/depth messages. Outputs camera-to-world poses, intrinsics, points and split & Valid image associations, dimensions, finite points and rotations\\
Training & Prepared scene, configuration and camera sequence. Outputs Gaussian parameters and metrics & Same initialization, exact ordered fitting indices, target RGB hashes and finite post-update parameters\\
Extraction & Final checkpoint and fitting cameras. Outputs raw PLY, OBJ and grid record & Matching checkpoint, crop, fusion cameras and settings for both arms\\
Cleanup & Raw PLY and extraction record. Outputs cleaned PLY and silhouette comparison & Input hash, finite loaded mesh, face and vertex counts and recorded filter operations\\
Editing & PLY, camera crop and fixed patch specification. Outputs Blender project and independent reopen record & Unit import transform, matching topology, saved document hash and evaluated displacement\\
\bottomrule
\end{tabularx}\par
\captionof{table}[Pipeline stages, outputs and integrity checks]{Stage interfaces and integrity checks.}\label{tab:interfaces}
\end{minipage}\end{center}

\subsection{Image preparation and camera estimation}\label{sec:acquisition}
\subsubsection{Public photograph route}
The primary indoor benchmark uses factor-two Mip-NeRF 360 images. The frozen outdoor Garden test uses factor-four images. Both routes read the prepared photographs and the distributed COLMAP models. The preparation script fetches the selected images and sparse model, copies the images into the prepared directory and rescales each camera to the actual width and height. The trainer uses data factor one on this prepared directory to avoid a second downsampling. The loader inverts world-to-camera extrinsics and normalizes camera centers and points using a common similarity followed by principal-axis alignment. It records a single transform for all points and cameras. The normalized camera distance bound defines the scene-scale factor used by the trainer. Scene scale governs thresholds and learning rates. It is distinct from the multiplicative normalization used to convert the matte depths back to meters.

The loader sorts image names and assigns every eighth image, beginning with index zero, to validation. The other views form the fitting set. Supplied SfM models use the benchmark images, so validation is a holdout from Gaussian RGB fitting. The camera model and sparse initializer may already reflect correspondences from those images. The project verifies readable dimensions, absence of exact duplicate groups and structurally valid rotations and points. It preserves low-texture views where the audit finds legitimate scene content, such as the blank wall in Room [\hyperlink{evE1}{E1}].

\subsubsection{Local RGB-D route}
The museum experiment begins with a RealSense D455 recording. Preparation preserves message identities, timestamps and file hashes. SDK-checked calibration and the replay convention transform depth points to the color camera. Frame provenance allows the split to be enforced at the depth-message level even when nearby images depict the same object. The selected matte segment has broader stable texture and better track support than the difficult exhibit and trajectory transitions. The whole-route repair experiments found unresolved intervals under independent checks, so the matched experiment uses a local coordinate system.

The local optimizer obtains tracks from fixed-calibration COLMAP features and sequential matches. Some feature tiles are withheld from fitting and used for independent correspondence checks. Track observations supply an initial robust metric scale. The joint bundle adjustment extends Equation~\ref{eq:ba} with metric camera-axis factors and fixes the first camera and the intrinsics. Factor weights account for range and local sample spread by an assumed uncertainty model. The bounded solve permits up to 45 evaluations. The fitting budget can terminate before numerical convergence, so the local consistency checks determine whether a submap can be used.

At approximately 10 fitting observations per second, the matte submap registered 92 of 92 fitting cameras. Its spatially withheld checks report median 3D separation 8.34 mm and median reprojection 0.468 pixels. These checks address consistency within the local route and declared support. The three reserved evaluation cameras are localized through RGB against the fixed map. Their depth identities are excluded from the fitting tracks, initializer and training-depth directory. Including those three poses gives 95 local camera records. The paired Gaussian test uses 21 selected fitting frames and three evaluation frames. It builds a common initializer of 90,000 evidence-ranked representatives. Supported points originate from fitting-depth samples and must have agreement from separated views with no observed free-space conflict [\hyperlink{evE2}{E2}].

\subsubsection{Video entry point and failure handling}
The auxiliary RGB-only MP4 pipeline extracts frames, runs CPU COLMAP, finds the largest model and checks registration before training. Its default gate requires at least 20 registered images, 100 sparse points and 80\% frame registration in that model. The tested user video failed the registration gate at both sampling rates. A high registration fraction is a minimum input check. Inspecting residuals, connected segments and independent revisits is still necessary before assigning physical meanings to uncertain regions. The public images-to-edit reproduction uses supplied SfM cameras and does not depend on the MP4 pipeline [\hyperlink{evE7}{E7}].

\subsection{Baseline Gaussian optimization}
Each prepared sparse point starts a Gaussian center. The initial scale is proportional to the root mean square distance to its three nearest neighbors, repeated along the three axes. The constant spherical harmonic term derives from the point color, and the higher-order coefficients begin at zero. Quaternions and opacity logits follow the trainer initialization with the fixed experiment seed. The same parameter tensors initialize the two arms in each comparison. A sorted tensor digest checks exact initial equality before any optimizer update.

At each update the trainer loads one fitting camera and differentiably rasterizes the projected scene. The active harmonic degree increases by one at each 1,000-update interval until degree three. The main supervised objective is
\begin{equation}
 L_{\mathrm{RGB}}=0.8\,\mathrm{mean}_{p,c}|C_{pc}-I_{pc}|
                         +0.2\bigl(1-\mathrm{SSIM}(C,I)\bigr).
\end{equation}
When a mask is supplied, the L1 term averages included pixels and the SSIM calculation zeros excluded positions on both sides. Adam optimizes each parameter family with its own learning rate. The center rate depends on scene scale and decays during training. Camera pose refinement is disabled in these runs, so the compared methods use the same camera model.

In the matte pair, the same confidence-weighted inverse-depth term is added to both arms with coefficient $0.01$. If $\widehat z_j$ is the bilinearly sampled expected depth, $z_j$ an aligned fitting depth and $q_j>0$ its confidence, the implemented residual is
\begin{equation}
 L_{\mathrm{depth}}=s_{\rm scene}
  \frac{\sum_{j\in\mathcal D}q_j\left|\operatorname{disp}(\widehat z_j)-1/z_j\right|}
       {\max\bigl(\sum_{j\in\mathcal D}q_j,10^{-8}\bigr)}.
\end{equation}
Here $\mathcal D$ contains valid positive finite sensor observations and weights, and the predicted disparity is zero for invalid or nonpositive predicted depth. The loss is zero on an empty support set. Surface-curvature and sensor-alpha penalties are disabled in the primary protocol. The curvature study varies its own coefficient in an earlier paired Room experiment.

B00 names the local trainer's \texttt{DefaultStrategy} with the study's refinement schedule. It uses signed projected-mean gradients at threshold $2\times10^{-4}$. The accumulated norm is averaged over projected observations. When the mean exceeds that threshold, the primitive is duplicated if its maximum scale is at most $0.01s_{\rm scene}$. Otherwise the parent is replaced with two randomly offset children. The stochastic offsets are sampled in local axes, transformed by the parent rotation, and scaled by the covariance axes. The children's scale is the parent scale divided by 1.6. There is no primitive ceiling for this arm. Both arms share opacity reset and the underlying pruning policy.

\subsection{A05-Abs budget-aware densification}\label{sec:allocation}
\subsubsection{Gradient gate}
For each projected occurrence of Gaussian $i$ in a render of width $W$ and height $H$, let $a_{iv}=\sum_p |b_{ivp}|$ be the renderer's existing component-wise absolute positional-gradient sum. The gradients correspond to the actual training objective, including the depth term where enabled. Let $V_i$ count the projected occurrences in the current accumulation interval. The trainer converts the projected position derivative to normalized screen coordinates. With one camera per batch the scalar statistic is
\begin{equation}
 g_i=\frac{\displaystyle\sum_{v:\,i\,\mathrm{projected}}
        \left\|\mathrm{diag}(W_v/2,H_v/2)\,a_{iv}\right\|_2}
        {\max(V_i,1)} .
\end{equation}
The generic gsplat strategy also scales by the number of cameras in the batch. Packed projections provide global IDs for this accumulation. A projected occurrence can count toward $V_i$ even if that Gaussian was absent from the top contributors for image scoring. With activated opacity $o_i$, the eligible set is
\begin{equation}
 \mathcal E=\{i:g_i>0.0008,\ o_i\geq0.005\}.
\end{equation}
The screen-radius branch and below-threshold fallback are disabled. The absolute-gradient threshold follows the example in the gsplat strategy documentation for this statistic. Raising the threshold changes eligibility in addition to switching the aggregation rule. The Room factorial measures both changes separately.

\subsubsection{Reconstruction-error attribution and edge evidence}
Starting at loop index 500, and every ten indices until the refinement stop, the strategy samples the current fitting image before backward. It does not compute a second RGB image. A CUDA contributor pass reuses the same projected means, conics, effective opacities and tile intersections. Each pixel reports up to four highest-weight contributors with their front-to-back weights and projected IDs. The observer remaps packed IDs to global primitive IDs and excludes ignored pixels, missing IDs and nonpositive or nonfinite weights. Attrition through compositing matters because a highly opaque primitive behind a visible surface can have little effect on the RGB residual.

Let $w_{ivp}$ be the compositing weight and $I_{vp}$ the target. The observer forms grayscale $J_{vp}$ by averaging the RGB channels. An absolute four-neighbor Laplacian with replicated borders provides edge strength $\ell_{vp}$. For every scored fitting image, define
\begin{align}
 e_{vp}&=\tfrac13\textstyle\sum_{c=1}^3 |C_{vpc}-I_{vpc}|,\\
 \ell_{vp}&=|J_{v,p+\hat x}+J_{v,p-\hat x}+J_{v,p+\hat y}+J_{v,p-\hat y}-4J_{vp}|.
\end{align}
If $\mathcal P_{iv}$ contains pixels where $i$ is in the reported contributor list, the strategy accumulates the sums
\begin{align}
 E_i &=\sum_{v}\sum_{p\in\mathcal P_{iv}}w_{ivp}e_{vp},&
 B_i &=\sum_{v}\sum_{p\in\mathcal P_{iv}}w_{ivp}\ell_{vp},&
 m_i &=\sum_v\mathbf 1[|\mathcal P_{iv}|>0].
\end{align}
Here $m_i$ counts contributing scored observations, once per image occurrence, and can count a repeated camera. The per-observation averages $E_i/\max(m_i,1)$ and $B_i/\max(m_i,1)$ preserve the effect of screen-space footprint area and contribution. They do not divide by the number of pixels or the sum of the weights. For example, if a pixel has average RGB error 0.20 and edge strength 0.40, a listed contributor with $w=0.30$ receives 0.06 error mass and 0.12 edge mass. Many pixels in the same scored image may add mass, while the observation count for that image remains one.

\subsubsection{Robust normalization and ranking}
Each score channel has a different numerical scale and distribution. For any vector $x$ across primitives, replace nonfinite and negative entries with zero to obtain $x_i^+$. Compute $m_x$ as the PyTorch median of positive $x_i^+$ over eligible parents. For even counts this takes the lower middle value. The normalization is
\begin{equation}
 {\cal N}_{\mathcal E}(x)_i =
 \begin{cases}
  \min(20,\,x_i^+/\max(m_x,10^{-12})), &\text{if an eligible positive value exists},\\
  0, &\text{otherwise}.
 \end{cases}
\end{equation}
This median scaling makes the score relative to the candidate population at the current event. It clips extreme normalized scores before combining them. There is no common global mean or variance. Normalized error and edge evidence are then discounted by $\eta_i=\sqrt{\min(1,m_i/3)}$. One scored observation contributes at about 0.577 of its normalized value, two at 0.816 and three or more at full weight. The gradient is unaffected by this discount. The final score is
\begin{equation}
 S_i={\cal N}_{\mathcal E}(g)_i+
 \eta_i\,{\cal N}_{\mathcal E}\bigl(E/\max(m,1)\bigr)_i+
 0.25\,\eta_i\,{\cal N}_{\mathcal E}\bigl(B/\max(m,1)\bigr)_i .
 \label{eq:ranking}
\end{equation}
The project uses equal weights for the normalized gradient and residual, with a reduced edge weight to limit the influence of texture and lighting contrast. These weights govern selection only. A Gaussian at a strong image edge can score highly without belonging to an actual depth boundary. Median references also shift when the eligibility rule changes, so the factorial contrasts measure the effect of the gradient rule through the complete selection mechanism.

\subsubsection{Event-indexed count ceiling}
Let $N_0$ be the initializer count and $C$ the configured cap. The schedule is defined on actual refinement events. For zero-based loop index $s$, it tests a strict start and stop, a periodic interval, and an optional pause after opacity reset. With the primary schedule the 54 events are $600,700,\ldots,5900$. Refinement at loop index 600 runs after the 601st update. Opacity reset triggers at index 3000, after that event's grow and prune. The reset at 6000 is skipped because the refinement callback returns at the stop index.

For event $k$ of $K$ scheduled events, define $u=k/K$. The schedule and the maximum number of additions are
\begin{align}
 C_k &=\min(C,\lfloor N_0+(C-N_0)(2u-u^2)\rfloor),\\
 b_k&=\max(0,C_k-N_k),\qquad
 a_k=\min(b_k,|\mathcal E_k|),
 \label{eq:count}
\end{align}
where $N_k$ is the live count before growth. The quadratic has a larger step early in refinement and a decreasing increment toward the cap. It permits earlier capacity allocation while making additions compete for the available slots. It neither forces growth to the cap nor uses the next event's pruning to fund current additions. Initial counts above $C$ raise an error.

Each selected parent has the same net cost of one primitive. If demand exceeds the slots, the strategy selects the $a_k$ highest scores with \texttt{topk}. Unselected parents stay in the model unless pruned. They are not queued and can compete again based on newly accumulated evidence. Exact score ties have the order returned by the backend. If demand is smaller than the slots, all eligible parents can split and the remainder is unused. If the count reaches the current target, no growth occurs. The next event recomputes headroom from its new target and live count.

For public Room at index 600, the cap schedule targets 134,180 from an initializer of 112,627. The 21,553 available slots exceed the 12,642 eligible parents. All eligible parents split, and 1,897 weak primitives are then pruned, leaving 123,372. At index 5900 the target reaches 700,000. There are 450,828 live primitives before growth. Eligibility limits additions to 1,958, and pruning removes 301, so the final count is 452,485. This example explains how the measured count can end far below the ceiling.

\subsubsection{Long-Axis Split and state transfer}
For a selected parent with longest scale $s_k$ and corresponding unit world axis $v_k=R_i\mathbf e_k$, the child centers and scales are
\begin{align}
 \mu_\pm&=\mu_i\pm 3d\,s_k v_k,&
 s'_k&=(1-d)s_k,&
 s'_j&=s_j\sqrt{1-d^2}\quad (j\ne k).
 \label{eq:las}
\end{align}
The construction in Figure~\ref{figure4} uses the parent's oriented $3\sigma$ extents. Along the original major axis, each child's $3\sigma$ half-width is $3(1-d)s_k$, so the outer endpoints of the pair coincide with the parent's original axial endpoints. At $d=0.45$, each center moves $1.35s_k$, each axial half-width is $1.65s_k$, and the axial intervals overlap by $0.60s_k$. The secondary scales retain about 89.3\% of their original values. This is a footprint construction using isocontour extents. It allows immediate changes to the rendered image and does not impose a physical continuity constraint.

\reportfigure[Child placement in Long-Axis Split]{0.89\linewidth}{figures/long_axis.png}
  {Long-Axis Split used by A05-Abs. The cross denotes the parent center. The child footprints are placed symmetrically along its longest rotated axis and shrink along all axes.}{figure4}

Both children copy the parent's quaternion and spherical harmonic coefficients. The optimizer stores $\log(s'_j)$ and $\operatorname{logit}(0.6\,o_i)$ for each child. Opacity reduction moderates overlapping contributions. It does not conserve pixel alpha, because transmittance and footprints change together. The $0.45$ offset and $0.6$ opacity multiplier use the empirical defaults of Deng et al. \cite{ref7}. Moving centers closer gives more projected overlap, and moving them farther can enlarge the discrepancy in the interior of the parent footprint. The archived A04 Room sweep found a small PSNR benefit for another setting in its selection run, with mixed results in confirmation. The primary candidate retains the defaults for comparison with the archived protocol.

The operator drops the selected parent rows, keeps the unselected rows and appends the positive-offset child rows followed by the negative-offset rows. The optimizer parameter-group reference is replaced for each attribute. The per-tensor Adam step counter persists and the surviving rows keep their first and second moments. Child moment rows are initialized to zero. Copying those moments from the parent would carry gradient history associated with its original location and covariance. Auxiliary per-primitive tensors are reindexed as the scene receives its split and removal hooks. Contribution and visibility statistics are zeroed for the next refinement window. Subsequent renders and optimizer steps use the new attributes and new per-row state.

\subsubsection{Refinement event and computational cost}
Figure~\ref{fig:event} summarizes one growth event. Table~\ref{tab:pseudocode} places it within the optimizer loop.
\begin{center}\begin{minipage}{\linewidth}
\centering\small
\setlength{\fboxsep}{5pt}
\begin{tabular}{@{}ccc@{}}
\fcolorbox{ReportBlue}{ReportLight}{\strut Average position gradients} &
$\longrightarrow$ &
\fcolorbox{ReportBlue}{ReportLight}{\strut Gate by gradient and opacity}\\
& & $\downarrow$\\
\fcolorbox{ReportBlue}{ReportLight}{\strut Normalize and rank eligible parents} &
$\longleftarrow$ &
\fcolorbox{ReportBlue}{ReportLight}{\strut Compute count headroom}\\
$\downarrow$ & & \\
\fcolorbox{ReportBlue}{ReportLight}{\strut Split up to the available slots} &
$\longrightarrow$ &
\fcolorbox{ReportBlue}{ReportLight}{\strut Prune, log and clear scores}
\end{tabular}\par
\captionof{figure}[Sequence of operations at a refinement event]{Sequence at a scheduled A05-Abs refinement event. The new live count includes both the split and the subsequent pruning.}\label{fig:event}
\end{minipage}\end{center}

\begin{center}\begin{minipage}{\linewidth}
\small
\begin{tabularx}{\linewidth}{@{}rY@{}}
\toprule
 & \textbf{Inputs} Prepared scene, ordered fitting-camera indices, initial Gaussians, Adam state and cap $C$.\\
\midrule
1 & Initialize counters and score masses. Reject an initializer that exceeds the cap.\\
2 & For each loop index $s$, load the next camera and rasterize the scene.\\
3 & \quad On every 10th index in the scoring interval, gather weighted contributor IDs and accumulate image evidence.\\
4 & \quad Backpropagate the fitting objective. Step the parameter optimizers and learning-rate schedules.\\
5 & \quad If $s<6000$, update the running projection-gradient norms and counts from that backward pass.\\
6 & \quad If $s$ is a scheduled refinement event, compute $C_k$, $a_k$ and the eligible set.\\
7 & \qquad If $a_k>0$, rank eligible parents by Equation~\ref{eq:ranking} and split the $a_k$ highest scores.\\
8 & \qquad Remove splats with $o_i<0.005$. After index 3000 also remove splats with maximum scale exceeding $0.1s_{\rm scene}$.\\
9 & \qquad Record the before, target, eligible, added, pruned and after counts. Clear the attribution and gradient statistics.\\
10 & \quad If $s=3000$, cap activated opacities at $0.01$ and reset their optimizer moments.\\
11 & \quad At requested evaluation indices, render the held-out views from the updated model.\\
12 & Save finite post-update parameters, camera order, metrics and growth history.\\
\bottomrule
\end{tabularx}\par
\captionof{table}[Training and refinement pseudocode]{A05-Abs training and refinement pseudocode. Loop indices are zero-based.}\label{tab:pseudocode}
\end{minipage}\end{center}

The additional observer uses $O(PK_p)$ storage for the reported IDs and weights at $P$ pixels and $K_p=4$. Selecting contributors still visits potential tile intersections, so computation depends on the number of overlapping projected Gaussians as well as image size. The three score channels require $O(N)$ running state. At a refinement event, the policy scans the live set, obtains a median per enabled channel on eligible candidates and performs a backend top-$a_k$ selection. An upper bound for a design using full sorting for those order statistics would be $O(N+|\mathcal E|\log|\mathcal E|)$. Actual GPU time also depends on memory movement when resizing the parameter and optimizer arrays. As $N$ grows, the cost of Adam and the rasterizer can dominate the extra score calculation. The runtime study measures the aggregate behavior rather than assuming a speedup from reducing $N$.

\subsubsection{Implementation map and design alternatives}
\begin{center}\begin{minipage}{\linewidth}
\small\begin{tabularx}{\linewidth}{@{}p{.25\linewidth}YY@{}}
\toprule
\textbf{Code path} & \textbf{B00 behavior} & \textbf{A05-Abs behavior}\\
\midrule
\texttt{simple\_trainer.py} & One camera per step, shared loss and optimizers & Same loop, with detached image-observer call and absolute-gradient render flag\\
\texttt{rendering.py} & Returns signed projected-mean gradients & Exposes absolute component gradients through the existing CUDA backward path\\
\texttt{strategy/default.py} & Gradient/size clone and random split, pruning, reset & Reused stage hooks, gradient accumulation, pruning and reset\\
\texttt{strategy/budget.py} & No allocation ranking or scheduled count ceiling & Candidate mask, image attribution, median scaling, support discount and budget history\\
\texttt{strategy/ops.py} & Duplicate or randomly split and migrate parameter rows & Symmetric split geometry and opacity, with the same optimizer/scene migration mechanism\\
\bottomrule
\end{tabularx}\par
\captionof{table}[Implementation of the baseline and candidate strategies]{Main code paths for the matched trainer.}\label{tab:code}
\end{minipage}\end{center}

Table~\ref{tab:code} identifies the files implementing each strategy. A constant cap would allow early events to use all spare capacity, potentially increasing cost for the rest of training. A rate limit based only on current count would be sensitive to opacity-reset pruning. The event schedule uses the initial count as a stable reference and current count as the allocation state. Gradient-only ranking is cheaper and provides a meaningful comparison, while target edges and residual attribution supply information about different aspects of the current error. The study disables each of those two image cues separately. The optional residual fallback was screened earlier against equal-count gradient selection and did not pass the practical criterion. The final candidate uses the ordinary gradient gate for every split.

The all-split policy avoids a clone at exactly the same center and leaves a common one-parent-to-two-children budget accounting rule. It is still possible for children to overlap substantially in the image. The earlier split-by-clone factorial distinguishes the operator bundle from the small-parent policy. Recovery-aware pruning and alternate update schedules from the related paper are outside the implemented policy. These choices keep the main trainer focused on how it grows Gaussians and preserve a common pruning rule for comparison.

\subsection{Gaussian depth rendering and surface extraction}
The extractor loads the final post-update checkpoint and reuses the trainer's camera parser. It renders at most 64 evenly spaced fitting cameras in the loader's fitting order. Both arms use the same camera identities. Resolution is reduced to $\lfloor W/4\rfloor$ by $\lfloor H/4\rfloor$, with the intrinsics scaled using the realized ratios. The render uses RGB and expected depth. Pixels require finite positive depth and accumulated alpha strictly greater than 0.7. Optional image masks exclude additional pixels. The integrator ignores observations beyond eight times the loader scene scale. Neither held-out views nor sensor depths are fused, including in the matte experiment.

For the uniform volume, compute the first and 99th initializer percentiles on each axis. Center the cube at the midpoint of those bounds and set its length $L$ to 1.2 times the largest span. With requested spacing $0.01s_{\rm scene}$ and maximum resolution 192, use
\begin{equation}
 R=\min(192,\lceil L/(0.01s_{\rm scene})\rceil),\qquad
 \Delta=L/R,\quad \tau=4\Delta.
\end{equation}
The cube and default grid settings are shared within a scene pair because both use the same initializer. The actual spacing can be coarser than requested. The bounded dense volume scales as $O(R^3)$ in voxel storage. It restricts unbounded backgrounds and keeps memory usage bounded. The crop, low-resolution expected depth and alpha cut all affect how far the reconstructed surface extends.

The implementation uses \texttt{UniformTSDFVolume} with RGB8 color and passes the world-to-camera inverse to the integrator. It exports PLY with vertex colors and normals and an OBJ suitable for standard editors. It rejects empty geometry, verifies finite vertices and reloads the PLY to check counts. The report saves the checkpoint digest, camera identities, occupancy support, grid and topology. A known-plane check justified using the uniform backend after the previous scalable backend returned an empty result on the same host. The fixed-truncation sensitivity study uses the identical uniform path with finer and coarser grids [\hyperlink{evE5}{E5}, \hyperlink{evE9}{E9}].

\subsection{Mesh cleanup and Blender verification}
The raw meshes use the same topology measurements before design cleanup. An edge incident to one face indicates an open boundary; incidence above two is an edge-level nonmanifold condition. Connected components use face adjacency through edges. These checks are computed on the loaded mesh. A watertight surface also requires closed boundaries and additional manifold and intersection conditions. Area or face fraction of the largest component helps distinguish many small fragments from a split in the bulk surface. A small component can be an intended object or an artifact.

The representative Room recipe removes degenerate and repeated triangular faces, including reversed duplicates. It then removes only connected components with fewer than eight faces and area less than a quarter of the squared voxel length. It compacts unused vertices and applies quadric decimation to approximately 65\% of the retained face count with a boundary penalty of 8. Quadric decimation ranks edge collapses by a local surface approximation error. The boundary penalty makes boundary changes more expensive without filling holes. A final duplicate/degeneracy pass and normal recomputation precede export. The cleanup record stores counts and hashes after each stage. The silhouette check uses the same cameras and crop boundaries on the raw and cleaned meshes.

The matched comparison selects a connected patch around a corresponding camera ray hit and displaces it toward the camera under a shared voxel-scaled recipe. The separate cleaned Room demonstration selects a connected tabletop patch around a fixed scene-coordinate seed, checks its height and upward normal, and lifts it using a cosine falloff. Blender stores the lift in a shape key at value one. A second headless process reopens the project, evaluates the mesh, verifies the selected displacement against the expected vector and checks that evaluating and rendering the file did not alter the saved document bytes. Photographs and meshes use the same chosen camera for visual comparison, and the Blender camera convention is obtained from the stored camera-to-world transform [\hyperlink{evE6}{E6}].

\subsection{Reproducibility and failure handling}
The trainer records configuration, initial tensors by digest, full camera order, metrics and the final post-update checkpoint. The legacy numbered checkpoint precedes its optimizer step. Extraction reads \texttt{final\_post\_update.pt}. Explicit camera replay addresses an earlier change in sample order after the first epoch when loader workers changed. Pair audits compare final RGB target pixels, extraction settings and camera lists. Data and checkpoint digests bind the results to their inputs.

Incomplete attempts remain available for diagnosis. The checks catch counts above the cap, mismatched depth identities, empty fusion, nonfinite tensors and unexpected Blender displacement. Appendix A lists the commands and expected artifacts.

\clearpage
\section{Implementation and Experimental Evaluation}
\subsection{Workstation, scenes and protocol}
The recorded workstation used Windows, an AMD Ryzen 5 3500, approximately 16 GiB system RAM and a GTX 1660 SUPER with 6 GiB device memory. The documented software environment includes Python 3.12.10, PyTorch 2.9.1+cu128, CUDA Toolkit 12.8.93 and a locally modified gsplat 1.6.0. Open3D 0.20.0 performed mesh cleanup and Blender 4.5.4 LTS performed the edits. These versions describe the recorded environment. Per-run configurations and hashes identify the experiments [\hyperlink{evE3}{E3}, \hyperlink{evE4}{E4}, \hyperlink{evE6}{E6}].\par
\begin{center}\begin{minipage}{\linewidth}
\centering\small\setlength{\tabcolsep}{4pt}
\begin{tabular*}{\linewidth}{@{\extracolsep{\fill}}lrrrl@{}}
\toprule
\textbf{Scene} & \textbf{Fitting / evaluation} & \textbf{RGB resolution} & \textbf{Initial primitives} & \textbf{Input geometry} \\
\midrule
Room & 272 / 39 & 1557 \ensuremath{\times} 1038 & 112,627 & Supplied SfM \\
Counter & 210 / 30 & 1558 \ensuremath{\times} 1038 & 155,767 & Supplied SfM \\
Kitchen & 244 / 35 & 1558 \ensuremath{\times} 1039 & 241,367 & Supplied SfM \\
Bonsai & 255 / 37 & 1559 \ensuremath{\times} 1039 & 206,613 & Supplied SfM \\
Museum matte & 21 / 3 & 834 \ensuremath{\times} 463 & 90,000 & Local RGB-D map \\
\midrule
Garden (follow-up) & 161 / 24 & 1297 \ensuremath{\times} 840 & 138,766 & Supplied SfM \\
\bottomrule
\end{tabular*}\par
\captionof{table}[Datasets, image splits and initial Gaussian counts]{Datasets and initialized primitive counts for the five primary pairs and the separate Garden follow-up. Public images are from Mip-NeRF 360. Museum frames are from the personal D455 recording.}\label{table2}
\end{minipage}\end{center}
Figure~\ref{figure5} shows one evaluation photograph from each of the five primary pairs. Each primary pair used 7,000 updates, seed 42, batch size one, packed rendering and spherical harmonics up to degree three. The refinement window ran at zero-based loop indices 600 through 5,900 in steps of 100. Opacity was reset at index 3,000. The candidate ceiling was 700,000 public and 180,000 matte. Public scenes used RGB supervision. Matte fitting frames span 6.008--13.817 seconds, with all three evaluation images inside that interval (Table~\ref{table2}). The 7,000-update comparison uses previously inspected development views, and its evidence is restricted to these scenes [\hyperlink{evE4}{E4}].\par
\reportfigure[Evaluation photographs from the five primary scene pairs]{0.95\linewidth}{figures/scene_photos.jpg}
  {Held-out photographs from the four indoor public scenes and the personal matte segment. The panels use the first evaluation target of each scene and are resized for display. Their resolutions and training conditions are given above.}{figure5}
\subsection{Evaluation criteria}
For normalized RGB, $\mathrm{PSNR}=-10\log_{10}(\mathrm{MSE})$ measures pixel fidelity in decibels. SSIM measures structural similarity \cite{ref10} and LPIPS compares learned image features \cite{ref11}. Higher PSNR and SSIM and lower LPIPS favor an arm. We use the final trainer averages over the same images. Raw surface evaluation counts edge-connected triangle components and edges with exactly one incident triangle. Depth discrepancy compares camera-axis depth on pixels observed by both methods. Reopening an edit tests a distinct operational outcome by checking whether a selected patch still moves in a design tool [\hyperlink{evE4}{E4}, \hyperlink{evE6}{E6}].\par
\clearpage
\subsection{Rendering quality and representation size}
\begin{center}\begin{minipage}{\linewidth}
\centering\small\setlength{\tabcolsep}{4pt}
\begin{tabular*}{\linewidth}{@{\extracolsep{\fill}}llrrrr@{}}
\toprule
\textbf{Scene} & \textbf{Arm} & \textbf{PSNR (dB)} & \textbf{SSIM} & \textbf{LPIPS} & \textbf{Gaussians} \\
\midrule
Room & B00 & 29.998 & 0.90168 & 0.20973 & 831,326 \\
Room & A05-Abs & 30.737 & 0.90915 & 0.19206 & 452,485 \\
\addlinespace[2pt]
Counter & B00 & 27.630 & 0.88807 & 0.19537 & 742,755 \\
Counter & A05-Abs & 27.840 & 0.88954 & 0.18812 & 381,228 \\
\addlinespace[2pt]
Kitchen & B00 & 29.916 & 0.91564 & 0.11661 & 1,082,362 \\
Kitchen & A05-Abs & 30.212 & 0.91519 & 0.11521 & 502,536 \\
\addlinespace[2pt]
Bonsai & B00 & 30.055 & 0.92395 & 0.16005 & 1,020,333 \\
Bonsai & A05-Abs & 30.391 & 0.92838 & 0.15284 & 463,678 \\
\addlinespace[2pt]
Matte & B00 & 28.872 & 0.88327 & 0.18209 & 485,156 \\
Matte & A05-Abs & 28.891 & 0.88560 & 0.18964 & 179,792 \\
\bottomrule
\end{tabular*}\par
\captionof{table}[Held-out RGB quality and final Gaussian counts]{Mean held-out RGB scores and final Gaussian counts at 7,000 updates. Higher PSNR and SSIM and lower LPIPS indicate better image agreement.}\label{table3}
\end{minipage}\end{center}
Figure~\ref{figure6} summarizes the paired count and PSNR differences in Table~\ref{table3}. All four public pairs favored A05-Abs in PSNR and LPIPS. Room gained 0.739 dB with 378,841 fewer primitives. Kitchen used 53.57\% fewer primitives and gained 0.297 dB, although SSIM fell by 0.00045. Across the public scenes, lower primitive count coincided with higher pixel fidelity. The much smaller matte PSNR change, 0.019 dB, accompanied worse LPIPS. Pooling those conditions would obscure the different source of depth information and the small personal evaluation set [\hyperlink{evE4}{E4}].\par
\reportfigure[Gaussian count reduction and PSNR change by scene]{0.95\linewidth}{figures/count_psnr.png}
  {Matched effect of A05-Abs relative to B00. Left, percentage reduction in final Gaussian count. Right, difference in mean held-out PSNR. Each row compares the same scene and training camera sequence. Matte uses depth supervision in both arms.}{figure6}
\clearpage
\subsection{Photographs and novel-view renders}
The shared validation targets allow direct inspection of what a scalar RGB score averages over. Figure~\ref{figure7} uses the identical photograph and the same source pixels in both models. Both models recover the plant silhouette. In the enlarged crop, the leaf edges and the cabinet seam behind the plant appear softer than in the photograph. The crop preserves the same field of view across the three panels, and the numerical interpretation uses all 39 views [\hyperlink{evE4}{E4}].\par
\reportfigure[Room photograph and novel-view renderings]{0.91\linewidth}{figures/room_rgb.jpg}
  {Public Room held-out view DSCF4891.JPG. Columns show the photograph, the default model and the budgeted model. The upper row locates the plant, and the lower row enlarges the same crop in each image. These are Gaussian RGB renderings before meshing.}{figure7}
Figure~\ref{figure8} asks a different question of the same learned representation. Ray casting a triangle mesh exposes missing surface coverage along the museum doorway and the display-case edges. Both methods recover broad panel outlines, while color and fine lettering blur. The candidate leaves a larger opening to the left of the doorway in the first reserved frame. A small gain in mean RGB PSNR can therefore coexist with a visible loss in an extracted surface [\hyperlink{evE4}{E4}].\par
\reportfigure[Museum doorway in the photograph and both raw meshes]{0.98\linewidth}{figures/matte_mesh.jpg}
  {Personal matte segment, held-out frame rgb\_0006487.png. Left to right, the recorded RGB image, B00 raw mesh and A05-Abs raw mesh rendered from the same pose. The bright doorway makes the broken surface boundary visible.}{figure8}
\clearpage
\subsection{Per-view variation and a Bonsai outlier}
The same-view comparisons in Table~\ref{table3} use the trainer's floating-point image metrics. For a separate distribution check, all 144 public and three matte final validation panels were decoded from their saved 8-bit RGB images. The target halves were verified byte for byte and against the pair record. The PSNR difference was then recomputed for each view at the saved precision. The direction of the scene means agrees with the main table, but individual gains vary substantially [\hyperlink{evE11}{E11}].

\begin{center}\begin{minipage}{\linewidth}
\small\begin{tabular*}{\linewidth}{@{\extracolsep{\fill}}lrrrr@{}}
\toprule
Scene & Improved / total & Median (dB) & Full mean (dB) & Without largest gain (dB)\\
\midrule
Room & 34 / 39 & $+0.306$ & $+0.740$ & $+0.582$\\
Counter & 26 / 30 & $+0.152$ & $+0.215$ & $+0.184$\\
Kitchen & 21 / 35 & $+0.048$ & $+0.293$ & $+0.240$\\
Bonsai & 20 / 37 & $+0.034$ & $+0.337$ & $-0.023$\\
Matte & 1 / 3 & $-0.064$ & $+0.017$ & $-0.386$\\
\bottomrule
\end{tabular*}\par
\captionof{table}[Per-view PSNR gains and outlier sensitivity]{Per-view PSNR differences from archived RGB panels. The sensitivity column recomputes the mean after dropping each scene's largest gain. The primary comparison keeps all views.}\label{tab:per-view}
\end{minipage}\end{center}

In Table~\ref{tab:per-view}, Room has the largest median gain, and Counter also improves on most targets. Gains are less consistent on Kitchen and Bonsai. One Bonsai view changes from 17.48 dB for B00 to 30.79 dB for A05-Abs. In Figure~\ref{fig:bonsai-extremes} the baseline has a broad dark cast, and the candidate reproduces more of the target brightness. Other views include a regression of 2.31 dB. Without the largest gain, the mean of the other 36 Bonsai views is slightly negative. This post hoc check shows that the positive Bonsai mean depends strongly on one view. Its cause and repeatability are unresolved. Similarly, the small matte mean gain masks two regressions among three views.

\reportfigure[Bonsai views with the largest PSNR gain and loss]{0.99\linewidth}{figures/bonsai_rgb_extremes.jpg}
  {Largest Bonsai PSNR gain and regression in archived final RGB panels. Columns show the same held-out photo and both Gaussian simulations. The selection illustrates extremes of the paired distribution.}{fig:bonsai-extremes}

\clearpage
\subsection{Computation and storage}
\begin{center}\begin{minipage}{\linewidth}
\centering\small\setlength{\tabcolsep}{4pt}
\begin{tabular*}{\linewidth}{@{\extracolsep{\fill}}llrrrrr@{}}
\toprule
\textbf{Scene} & \textbf{Arm} & \textbf{Train (s)} & \textbf{Extract (s)} & \textbf{Peak (GiB)} & \textbf{Render (fps)} & \textbf{File (MiB)} \\
\midrule
Room & B00 & 424.6 & 8.62 & 1.300 & 75.4 & 187.1 \\
Room & A05-Abs & 394.8 & 8.16 & 0.971 & 111.0 & 101.8 \\
\addlinespace[2pt]
Counter & B00 & 441.7 & 13.28 & 1.185 & 57.8 & 167.2 \\
Counter & A05-Abs & 432.8 & 10.72 & 0.945 & 92.8 & 85.8 \\
\addlinespace[2pt]
Kitchen & B00 & 560.2 & 16.27 & 1.707 & 44.6 & 243.6 \\
Kitchen & A05-Abs & 469.8 & 12.96 & 1.077 & 67.4 & 113.1 \\
\addlinespace[2pt]
Bonsai & B00 & 449.6 & 16.65 & 1.588 & 58.3 & 229.6 \\
Bonsai & A05-Abs & 398.8 & 16.83 & 1.006 & 118.5 & 104.4 \\
\addlinespace[2pt]
Matte & B00 & 214.5 & 4.40 & 0.756 & 146.1 & 109.2 \\
Matte & A05-Abs & 190.6 & 4.51 & 0.342 & 249.7 & 40.5 \\
\bottomrule
\end{tabular*}\par
\captionof{table}[Computation, memory, rendering throughput and storage]{Observed runtime, peak tensor allocation, evaluation render throughput and checkpoint size. Each row describes one run.}\label{table4}
\end{minipage}\end{center}
In Table~\ref{table4}, extraction includes loader and checkpoint setup, rendering, fusion, export and internal checks. Training time uses the trainer loop through the last update, including the initial evaluation. Peak memory is the trainer's maximum PyTorch tensor allocation. Render rate is the reciprocal of mean CUDA-synchronized render-call time on the final evaluation images at their listed resolutions. File size is the byte length of the Gaussian checkpoint passed to extraction, divided by $2^{20}$. These quantities describe distinct parts of the system [\hyperlink{evE4}{E4}].\par
The Room checkpoint fell from 187.1 to 101.8 MiB, and the timed render call rose from 75.4 to 111.0 frames/s. The latter rates exclude image loading, display, encoding and viewer interaction. A05-Abs reduced peak tensor allocation on every scene. Extraction time did not follow count as closely: the same grid still had to be fused and meshed, and Bonsai took 16.83 s with A05-Abs versus 16.65 s with B00. The archived Bonsai baseline used four data workers and the candidate used zero, so its timing also includes a loader difference [\hyperlink{evE4}{E4}]. Figure~\ref{figure9} explains why Room finished far below the scheduled ceiling. At the last refinement event, 1,958 eligible parents split and 301 primitives were pruned, giving 452,485 live Gaussians [\hyperlink{evE3}{E3}, \hyperlink{evE4}{E4}].\par
\reportfigure[Room Gaussian count and scheduled ceiling during refinement]{0.79\linewidth}{figures/growth.png}
  {Public Room, seed 42. The dashed curve shows the scheduled growth ceiling. The solid curve shows the count after each recorded A05-Abs refinement event, including pruning. The gap reflects the eligibility-limited allocation in this run.}{figure9}
\clearpage
\subsection{Controlled Room study}\label{sec:room-study}
The follow-up replays 7,000 fitting updates at seed 42 with the current code, the archived fitting camera order and the same 39 evaluation targets. Both primary arms use zero loader workers. All roles share the same fitting objective, initialization and camera order. Every arm is audited against the known initializer and the target RGB hashes. This study is a new measurement. The original five-scene comparisons above retain their own endpoint values, and comparisons in this section use only runs from the same replay.

The Room protocol crosses gradient aggregation (signed or absolute) with eligibility threshold ($2\times10^{-4}$ or $8\times10^{-4}$) under the same cap, ranking and split operator. It also disables RGB-error attribution and edge strength one at a time. Three cap settings measure the realized size and quality at the default absolute threshold. The cap sweep showed that the 700,000 ceiling did not restrict A05-Abs selection on Room. Two additional runs remove each ranking term at the binding 400,000 cap, where the remaining scores decide which eligible parents receive the available slots.

\subsubsection{Matched updates and elapsed training time}
The main pair also evaluates after updates 1, 1,000, 2,000, \ldots, 7,000 and at target elapsed times 90, 180 and 270 seconds. Time starts at the training loop after loading and initialization. CUDA synchronization measures completed work. Deadline checks occur at 10-update intervals, and elapsed time excludes evaluation. Each sample is taken just after crossing its time target. Table~\ref{tab:matched-time} reports the actual times and counts. The same elapsed budget can correspond to a different number of optimizer updates. Figure~\ref{fig:room-progress} plots the trajectories.

\begin{center}\begin{minipage}{\linewidth}
\small\begin{tabular*}{\linewidth}{@{\extracolsep{\fill}}rlrrrrrr@{}}
\toprule
Target (s) & Arm & Actual (s) & Updates & PSNR (dB) & SSIM & LPIPS & Gaussians\\
\midrule
\inputrows{evidence/room_time_rows.tex}
\bottomrule
\end{tabular*}\par
\captionof{table}[Room results at matched elapsed training times]{Public Room at shared elapsed training targets. Time excludes interleaved evaluation.}\label{tab:matched-time}
\end{minipage}\end{center}

\reportfigure[Room learning curves by update count and elapsed time]{0.98\linewidth}{figures/room_progress.pdf}
  {Room learning trajectories. The update panel uses fixed update milestones. The time panel uses the actual seconds measured near the three common deadlines. Both evaluate the same 39 development views. The update panel omits the nearly identical first-update scores to show the later differences.}{fig:room-progress}
At the 90, 180 and 270 second targets, the candidate led by 0.099, 0.261 and 0.513 dB, respectively, with fewer Gaussians at each sample. At 90 seconds its SSIM and LPIPS were slightly worse. At 270 seconds it completed 5,290 updates versus 5,390 for B00, while using 432,794 versus 781,109 Gaussians and scoring 30.236 versus 29.723 dB. The candidate accumulated splats more quickly through the first 1,000 updates. Between 1,000 and 2,000 updates, the baseline count overtook it, while the candidate maintained a PSNR lead. At the sampled deadlines, the PSNR gain and count reduction occurred despite fewer completed updates.

The final replay scores were 30.757 dB at 452,409 Gaussians for A05-Abs and 30.125 dB at 834,493 for B00. This is a 0.632 dB difference and a 45.8\% count reduction. Both scores continued to rise between the 6,000- and 7,000-update probes with counts unchanged, reflecting attribute optimization after structural refinement stopped. The paired replay reproduced the direction of the archived Room result. Exact scores and counts differed slightly from the archive, so the replay and archived values are reported separately [\hyperlink{evE8}{E8}].


\subsubsection{Gradient aggregation, threshold and ranking}
For a metric $M$, let $M_{a,h}$ denote the absolute-gradient result at the higher threshold, with $s$ denoting signed and $l$ the lower threshold. The effect of changing aggregation at fixed high threshold is $M_{a,h}-M_{s,h}$. The threshold effect within absolute gradients is $M_{a,h}-M_{a,l}$. The interaction is $(M_{a,h}-M_{s,h})-(M_{a,l}-M_{s,l})$. Table~\ref{tab:room-components} gives the four single-seed arms. Table~\ref{tab:tight-ranking} repeats the ranking contrasts at a binding ceiling. Counts and appearance jointly reflect the accumulated changes to eligibility and parent ranking.

\begin{center}\begin{minipage}{\linewidth}
\small\begin{tabular*}{\linewidth}{@{\extracolsep{\fill}}llrrrrr@{}}
\toprule
Policy & Aggregation & Threshold & PSNR (dB) & SSIM & LPIPS & Gaussians\\
\midrule
\inputrows{evidence/room_component_rows.tex}
\bottomrule
\end{tabular*}\par
\captionof{table}[Room gradient and ranking ablations at a 700,000 cap]{Room gradient-by-threshold factorial and image-ranking ablations at a 700,000-Gaussian cap. All rows use Long-Axis Split and 7,000 updates.}\label{tab:room-components}
\end{minipage}\end{center}
\begin{center}\begin{minipage}{\linewidth}
\small\begin{tabular*}{\linewidth}{@{\extracolsep{\fill}}lrrrrr@{}}
\toprule
Ranking policy & Binding events & PSNR (dB) & SSIM & LPIPS & Gaussians\\
\midrule
\inputrows{evidence/room_ranking_rows.tex}
\bottomrule
\end{tabular*}\par
\captionof{table}[Ranking ablations at a binding 400,000 cap]{Ranking ablations at the binding 400,000-Gaussian cap. All arms use absolute gradients and threshold $8\times10^{-4}$.}\label{tab:tight-ranking}
\end{minipage}\end{center}
Absolute aggregation raised PSNR by 0.487 dB at the low threshold and by 1.089 dB at the high threshold, an interaction of about 0.602 dB. It also increased count at each fixed threshold. The absolute low-threshold arm reached 31.084 dB and 699,645 Gaussians, with competition for slots at all 54 events. The selected high-threshold arm gave up 0.327 dB and saved 247,236 primitives (35.3\%). It still improved on B00's 30.125 dB by 0.632 dB with 45.8\% fewer primitives. The high-threshold signed arm grew to only 143,086 Gaussians and fell below B00 on all three RGB metrics. These contrasts measure the gradient eligibility policy under fixed split and ranking settings. The low threshold also activates the cap, so the threshold contrast includes that downstream effect.

The earlier A04 factorial supports a contribution from Long-Axis Split and replacing clones with splits under signed gradients (Table~\ref{tab:split-ablation}). An equivalent factorial under A05-Abs would measure whether the split effects change with absolute gradients [\hyperlink{evE3}{E3}].

At 700,000, every eligible parent split in all three absolute/high-threshold ranking arms. At 400,000, competition occurred in 32 or 33 events. Removing the RGB-error cue changed PSNR by $-0.018$ dB and removing the edge cue changed it by $+0.033$ dB, at nearly equal counts. Both removals slightly increased LPIPS. The full configuration remains the fixed reference because it is the design carried forward from development and has the lowest LPIPS in this binding-cap trio. The edge-free variant has the highest PSNR here. A single Room seed and opposing metric directions leave the general value of the cues unresolved [\hyperlink{evE8}{E8}].


\subsubsection{Gaussian cap and realized count}\label{sec:room-caps}
Sweeping the cap tests how the same strategy behaves under different count ceilings. Each setting uses the same gradient definition and threshold, ranking, split operation and 7,000 training updates. The cap changes the event targets throughout training. It can alter which parents survive and compete later, so two similar endpoints can still arise through different trajectories. The baseline contributes one unrestricted endpoint for context.

\begin{center}\begin{minipage}{\linewidth}
\small\begin{tabular*}{\linewidth}{@{\extracolsep{\fill}}lrrrrr@{}}
\toprule
Arm & Cap & Gaussians & PSNR (dB) & SSIM & LPIPS\\
\midrule
\inputrows{evidence/room_cap_rows.tex}
\bottomrule
\end{tabular*}\par
\captionof{table}[Room cap sweep at equal update count]{Room cap sweep at equal update count. The final count is the measured model size.}\label{tab:room-caps}
\end{minipage}\end{center}

\reportfigure[Room PSNR against realized Gaussian count]{0.92\linewidth}{figures/room_cap_tradeoff.pdf}
  {Room quality against realized Gaussian count at 7,000 updates. The three budgeted points highlight the tighter cap, the selected high threshold and the low threshold. Table~\ref{tab:room-caps} also gives the 1,000,000-ceiling control, which overlaps the selected operating point at this scale.}{fig:cap-tradeoff}
At the selected threshold, both the 700,000 and 1,000,000 schedules admitted every eligible parent and finished near 450,000 Gaussians. Top-$k$ ranking and cap enforcement excluded zero parents in the selected 700,000 run. Its gain over B00 measures the bundled changes in eligibility and splitting. The exact contribution of each change remains unresolved. The ceiling limits growth where eligibility exceeds capacity.

At 400,000, the guardrail constrained 33 of 54 events. It saved 52,590 primitives relative to the 700,000-cap arm at a cost of 0.051 dB. It still led B00 by 0.581 dB with 52.1\% fewer primitives. Figure~\ref{fig:cap-tradeoff} shows the resulting quality--representation-size choices. A larger model bought higher PSNR at the low absolute threshold, and a stricter cap bought a smaller model at a measurable quality cost [\hyperlink{evE8}{E8}].


\clearpage
\subsection{Room seed replication and outdoor evaluation}\label{sec:followup}
The protocol fixed the policy, splits, seeds and outcomes before training. Room seeds 43 and 44 extend the seed-42 replay on the same factor-two split. Garden, outside method selection, uses factor-four images, supplied SfM, seed 42 and 24 reserved views. Each pair shared its initializer, 7,000 fitting indices, target RGB pixels and zero-worker loader. Table~\ref{tab:followup} gives the endpoints. Figure~\ref{fig:paired-followup} plots the paired differences [\hyperlink{evE12}{E12}].

\begin{center}\begin{minipage}{\linewidth}
\small\begin{tabular*}{\linewidth}{@{\extracolsep{\fill}}lrlrrrrr@{}}
\toprule
Scene & Seed & Arm & PSNR (dB) & SSIM & LPIPS & Gaussians & Binding events\\
\midrule
\inputrows{evidence/followup_rows.tex}
\bottomrule
\end{tabular*}\par
\captionof{table}[Paired Room seed replications and outdoor Garden results]{Paired final RGB scores, Gaussian counts and growth events with more eligible parents than slots. The ceiling applies to A05-Abs.}\label{tab:followup}
\end{minipage}\end{center}

The mean paired Room PSNR difference was \(+0.674\pm0.042\) dB (mean $\pm$ sample SD across three seeds). Gaussian count fell by 45.6\% on average, with a sample SD of 0.3 percentage points. These summaries describe repeatability on Room under this protocol.

\reportfigure[Paired count reductions and PSNR gains on Room and Garden]{0.83\linewidth}{figures/paired_followup.pdf}
  {A05-Abs relative to B00 at 7,000 updates. Each row pairs the same scene, seed and fitting order. Garden has one outdoor pair at factor four.}{fig:paired-followup}
All three Room pairs improved on SSIM and LPIPS. Every eligible parent split at every seed. These replications assess the complete policy.

Garden gained 0.365 dB with 75.6\% fewer Gaussians. Here the ceiling constrained all 54 events. In the saved 8-bit panels, 23 of 24 views gained PSNR, with a median gain of 0.318 dB. Removing the largest gain still left a positive mean of 0.335 dB. This outdoor result covers one seed within the same benchmark family [\hyperlink{evE12}{E12}].


\clearpage
\subsection{Earlier allocation and surface-prior controls}
An earlier study compared the budgeted signed-gradient predecessor A04 with B00 across three paired seeds on all four public scenes. Mean count fell by 24.25--35.27\% and mean PSNR rose on each scene. The paired intervals for Kitchen and Bonsai included zero. A second study compared A05-Abs directly with A04 on Room, Counter and Kitchen at seeds 42 and 43. Counts fell and PSNR rose in all six comparisons. Five passed the predeclared criteria of at least +0.05 dB PSNR, at most 0.001 SSIM loss and at most 0.001 LPIPS increase. Counter at seed 43 gained 0.04191 dB. As Figure~\ref{figure10} shows, the two Counter estimates straddle the threshold [\hyperlink{evE3}{E3}].\par
\reportfigure[Absolute-gradient comparison with the signed-gradient predecessor]{0.84\linewidth}{figures/ablation.png}
  {A05-Abs versus its budgeted predecessor A04. Each symbol is one paired difference at 7,000 updates. The dashed line marks the PSNR part of the acceptance criterion. Counter at seed 43 falls below it. Measurements come from the paired absolute-gradient study.}{figure10}
On Room, adding an edge-aware log-depth curvature penalty at weight 0.01 from update 3,000 lowered the reduced model's common-support curvature by 3.06\%. Its raw mesh nevertheless rose from 1,162 to 1,216 components and from 16,807 to 17,700 boundary edges. Table~\ref{table5} shows the same prior applied to the baseline. The opposite changes in component count across the two model families motivated retaining the no-prior candidate. A smoothness penalty can reduce a depth derivative while leaving the continuity of the extracted surface unresolved [\hyperlink{evE5}{E5}].\par
\begin{center}\begin{minipage}{\linewidth}
\centering\small\setlength{\tabcolsep}{4pt}
\begin{tabular*}{\linewidth}{@{\extracolsep{\fill}}lrrrrr@{}}
\toprule
\textbf{Room model} & \textbf{Prior weight} & \textbf{PSNR (dB)} & \textbf{LPIPS} & \textbf{Gaussians} & \textbf{Components} \\
\midrule
B00 & 0 & 29.99812 & 0.209726 & 831,326 & 1,241 \\
B00 & 0.01 & 30.05953 & 0.210307 & 827,764 & 1,131 \\
A05-Abs & 0 & 30.73701 & 0.192058 & 452,485 & 1,162 \\
A05-Abs & 0.01 & 30.73650 & 0.191455 & 450,398 & 1,216 \\
\bottomrule
\end{tabular*}\par
\captionof{table}[Effect of the depth-curvature penalty on Room]{Room ablation at seed 42. The curvature penalty reduces fragmentation on B00 and increases it on the reduced model.}\label{table5}
\end{minipage}\end{center}
\clearpage
\subsection{Earlier split controls and hyperparameter sensitivity}
The earlier A04 campaign varied the split operator and the treatment of small parents in a $2\times2$ factorial. It held the budget, ranking and camera order fixed within each of three Room seeds. The random operator uses the original offset sampling, scale reduction and opacity convention. The Long-Axis Split operator changes placement, scale and opacity as one bundle. The small-parent policy either clones those parents or applies the same split operator used for large ones.

\begin{center}\begin{minipage}{\linewidth}
\small\begin{tabular*}{\linewidth}{@{\extracolsep{\fill}}llrrrr@{}}
\toprule
Operator & Small parents & PSNR (dB) & SSIM & LPIPS & Gaussians\\
\midrule
Random & Clone & $30.129\pm0.090$ & 0.90202 & 0.21174 & 699,202\\
Random & Split & $30.292\pm0.120$ & 0.90360 & 0.20486 & 699,739\\
Long axis & Clone & $30.425\pm0.076$ & 0.90332 & 0.21023 & 699,226\\
Long axis & Split & $30.633\pm0.077$ & 0.90552 & 0.20344 & 598,107\\
\bottomrule
\end{tabular*}\par
\captionof{table}[Split operator and small-parent policy across three seeds]{Room A04 split operator by small-parent policy, paired seeds 42--44. Shown are mean and sample SD for PSNR and means for the other outcomes.}\label{tab:split-ablation}
\end{minipage}\end{center}

With cloning retained, Long-Axis Split gained $0.296\pm0.082$ dB over random splitting. With all selected parents split, the corresponding gain was $0.341\pm0.132$ dB. Replacing clones under the long-axis operator gained $0.208\pm0.022$ dB and lowered the final count. These paired effects support the combined operator and all-split design on this Room protocol. The operator changes three child attributes together. The experiment measures that bundle on A04 with signed gradients [\hyperlink{evE3}{E3}].

\begin{center}\begin{minipage}{\linewidth}
\small\begin{tabular*}{\linewidth}{@{\extracolsep{\fill}}rrrrrr@{}}
\toprule
Distance $d$ & Opacity multiplier & PSNR (dB) & SSIM & LPIPS & Gaussians\\
\midrule
0.30 & 0.40 & 30.6196 & 0.90462 & 0.20598 & 525,570\\
0.30 & 0.80 & 30.5459 & 0.90360 & 0.20538 & 568,339\\
0.50 & 0.40 & 30.6061 & 0.90539 & 0.20281 & 607,499\\
0.50 & 0.80 & 30.4872 & 0.90345 & 0.20352 & 641,270\\
\textbf{0.45} & \textbf{0.60} & 30.6007 & 0.90520 & 0.20279 & 597,975\\
\bottomrule
\end{tabular*}\par
\captionof{table}[Sensitivity to split distance and opacity multiplier]{Earlier A04 Room seed-42 sensitivity to the split distance and opacity multiplier. The reference row is the setting retained for A05-Abs.}\label{tab:offset-ablation}
\end{minipage}\end{center}

In Table~\ref{tab:offset-ablation}, the $(0.30,0.40)$ setting led PSNR in the Room selection run by about 0.019 dB, with fewer primitives and higher LPIPS. Its confirmation PSNR changes relative to the reference were $+0.035$ and $-0.017$ dB on the two remaining Room seeds and $+0.027$ dB on Counter seed 42. LPIPS increased in each confirmation case. The split inherits the $0.45$ distance and $0.60$ opacity defaults supported as empirical trade-offs in the source literature. The local sensitivity data show that another setting can move count or PSNR favorably without jointly improving the image metrics across the tested repeats. The reference remains a documented operating point [\hyperlink{evE3}{E3}].

The earlier signed-gradient cap sweep on Room produced 399,842, 597,975 and 604,705 Gaussians at caps of 400,000, 700,000 and 1,000,000. Their PSNR values were 30.546, 30.601 and 30.605 dB. This behavior illustrates the distinction between an imposed ceiling and the count that eligibility produces. The new A05-Abs sweep in Section~\ref{sec:room-caps} tests whether the same pattern persists after switching gradient mode and threshold [\hyperlink{evE3}{E3}, \hyperlink{evE8}{E8}].

\subsection{Interpreting geometry and failure cases}
Three causes can produce a missing or displaced surface. Camera errors move the predicted locations of observations between views. Representation errors change the expected depths before meshing. Extraction choices change which depth estimates contribute and how they become triangles. The public pair holds camera input and extractor settings fixed, so a paired difference reflects how the two learned models interact with that extractor. The public comparisons cannot identify absolute error in the shared input camera models. The matte experiment evaluates Gaussian depths before fusion and mesh depths afterward. The grid sweep measures one extraction variable while the checkpoints are fixed.

Table~\ref{tab:failure-cases} connects the inspected defects to the surfaces and views where they occurred. Narrow supports disappear in mesh renders of Room and thin tread detail fuses in Kitchen even when the corresponding RGB view is recognizable. In Bonsai the model reproduces the color pattern of the blossoms while the mesh joins many blossoms into a few surfaces. The matte doorway has a larger opening in the candidate mesh, and the bright boundary is an uncertain transition for opacity and depth. These examples show what the topology counts alone cannot diagnose. Each is selected for inspection and does not estimate the frequency of that failure mode across all pixels or scenes.

\begin{center}\begin{minipage}{\linewidth}
\small\begin{tabularx}{\linewidth}{@{}p{.23\linewidth}YY@{}}
\toprule
\textbf{Case} & \textbf{Observation} & \textbf{Diagnostic scope}\\
\midrule
Room tabletop and supports & Broad top survives. Thin supports remain incomplete before and after simplification & The tabletop lift samples the selected patch\\
Kitchen toy bulldozer & Outline survives. Adjacent tread details merge at common extraction settings & The RGB image can represent texture in the merged region\\
Bonsai and bicycle & Blossom clusters join. Bicycle details break apart & Counts mix separate objects with fragments\\
Matte doorway and glass edge & Candidate opening grows while stable-depth coverage increases overall & Pooled coverage can increase on one surface and decrease on another\\
Museum trajectory & Independent validation leaves many route intervals unresolved & Matte evidence applies to one local submap\\
\bottomrule
\end{tabularx}\par
\captionof{table}[Observed surface and camera failure cases]{Observed failure modes and the scope of their inspection.}\label{tab:failure-cases}
\end{minipage}\end{center}

\clearpage
\subsection{Raw surface quality}
\begin{center}\begin{minipage}{\linewidth}
\centering\small\setlength{\tabcolsep}{4pt}
\begin{tabular*}{\linewidth}{@{\extracolsep{\fill}}llrrr@{}}
\toprule
\textbf{Scene} & \textbf{Arm} & \textbf{Triangles} & \textbf{Components} & \textbf{Boundary edges} \\
\midrule
Room & B00 & 119,613 & 1,241 & 18,165 \\
Room & A05-Abs & 119,987 & 1,162 & 16,807 \\
\addlinespace[2pt]
Counter & B00 & 158,759 & 2,606 & 33,559 \\
Counter & A05-Abs & 163,497 & 2,443 & 33,701 \\
\addlinespace[2pt]
Kitchen & B00 & 70,332 & 661 & 6,446 \\
Kitchen & A05-Abs & 82,518 & 944 & 9,148 \\
\addlinespace[2pt]
Bonsai & B00 & 322,809 & 5,364 & 66,045 \\
Bonsai & A05-Abs & 320,840 & 4,859 & 64,084 \\
\addlinespace[2pt]
Matte & B00 & 80,952 & 973 & 13,300 \\
Matte & A05-Abs & 99,107 & 1,448 & 21,023 \\
\bottomrule
\end{tabular*}\par
\captionof{table}[Raw mesh topology at matched extraction settings]{Raw TSDF mesh topology at matched per-scene extraction settings. Component counts refer to edge-connected triangles.}\label{table6}
\end{minipage}\end{center}
Across the ten raw meshes in Table~\ref{table6}, no edge had more than two incident triangles. Their open boundaries rule out watertightness. A separate Open3D readback also failed vertex manifoldness on every raw mesh and the cleaned Room mesh. All passed the orientability check. Their largest edge-connected components covered 82.8--95.6\% of raw mesh area. Room and Bonsai had fewer components and boundary edges with A05-Abs, while Kitchen and matte had more of both. Counter had fewer components and slightly more boundary edges. Figure~\ref{figure11} shows the relative changes. The counts include intended objects as well as fragments [\hyperlink{evE4}{E4}, \hyperlink{evE11}{E11}].\par
\reportfigure[Change in raw mesh topology across five scenes]{0.95\linewidth}{figures/topology.png}
  {Change in raw mesh topology under identical per-scene extraction settings. Bars show percentage change from B00. An increase in boundary edges indicates more exposed triangle borders, whose location and cause require visual inspection.}{figure11}
On Kitchen, both meshes preserve the outline of the toy bulldozer but merge tread detail. Bonsai blossoms merge into larger surfaces, and background bicycle details break apart. On Room the largest structures remain recognizable even where thin supports fail. These observations are consistent with averaging depths at mixed pixels and sampling them on a finite grid. Section~\ref{sec:grid-results} probes grid sensitivity while holding the two matte checkpoints fixed [\hyperlink{evE4}{E4}].\par
\clearpage
\subsection{Measured-depth consistency in the museum}
Three reserved depth messages supplied 977,442 valid depth pixels between 0.3 and 6 m. A local 3 \ensuremath{\times} 3 range filter retained 918,331 stable pixels out of 1,158,426 total image pixels. Each method's coverage is reported against this stable set. Discrepancies are pooled over the exact pixels shared by the pair: 916,892 for Gaussian depth and 881,411 for raw mesh first hits. Table~\ref{table7} summarizes the resulting discrepancies. Negative signed discrepancy means the prediction is closer to the camera than the aligned sensor measurement [\hyperlink{evE2}{E2}, \hyperlink{evE4}{E4}].\par
\begin{center}\begin{minipage}{\linewidth}
\centering\small\setlength{\tabcolsep}{4pt}
\begin{tabular*}{\linewidth}{@{\extracolsep{\fill}}llrrrr@{}}
\toprule
\textbf{Prediction} & \textbf{Arm} & \textbf{Coverage} & \textbf{Signed med. (cm)} & \textbf{Abs. med. (cm)} & \textbf{Abs. P90 (cm)} \\
\midrule
Gaussian expected depth & B00 & 99.95\% & \ensuremath{-}6.64 & 7.37 & 21.62 \\
Gaussian expected depth & A05-Abs & 99.88\% & \ensuremath{-}7.33 & 7.97 & 21.25 \\
Raw mesh first hit & B00 & 97.30\% & \ensuremath{-}8.32 & 8.38 & 20.20 \\
Raw mesh first hit & A05-Abs & 97.82\% & \ensuremath{-}9.34 & 9.47 & 20.28 \\
\bottomrule
\end{tabular*}\par
\captionof{table}[Gaussian and mesh depth against the matte sensor observations]{Held-out sensor agreement in the matte segment. Coverage uses all stable pixels. Discrepancies use the common support within each output pair.}\label{table7}
\end{minipage}\end{center}
The candidate mesh covered 0.52 percentage points more stable depth, yet its median absolute discrepancy increased by 1.09 cm. The median worsened on each reserved view. Gaussian expected depth and mesh first-hit depth answer different questions near transparent layers and gaps, which is why both appear in the table. Neither provides an independent physical reference: the scale, calibration and camera estimates derive from the same capture system [\hyperlink{evE4}{E4}].\par
A whole-route diagnostic screened 317 eligible museum intervals. Only 132 passed, and none of the 185 failed intervals passed on retry. A source-following holdout later distinguished a passing matte control from a difficult target that still failed reprojection. This justified restricting the paired depth experiment to local matte surfaces. Reconstruction of the complete museum remains deferred. An older comparison on a separate RealSense sequence improved PSNR from 18.887 to 20.300 dB on the same 79 targets when sampling, matching and pose recovery changed together. It does not identify a frame-rate-only effect [\hyperlink{evE2}{E2}, \hyperlink{evE7}{E7}].\par
\clearpage
\subsection{Fixed-checkpoint TSDF grid sensitivity}\label{sec:grid-results}
The matched reconstruction table uses one TSDF grid per scene. A separate matte experiment retains both final Gaussian checkpoints and compares grid resolutions 128, 192 and 256 over the same crop. It holds the 21 fitting camera samples, their rendering resolution, alpha threshold and absolute truncation length fixed. Resolution 192 uses the archived extraction for each arm. The other four meshes are freshly extracted. The actual normalized voxel lengths are approximately 0.0647, 0.0431 and 0.0324, respectively. The truncation distance stays at 0.1726 normalized units.

For each reserved sensor frame, cast the same calibrated rays into all six meshes. Also render expected depth from both fixed Gaussian checkpoints. Compare discrepancies on the intersection of stable sensor pixels and valid predictions from all eight outputs. This is a stricter support than either original pairwise comparison. Table~\ref{tab:grid-depth} reports each output's coverage over all stable pixels. Table~\ref{tab:grid-topology} gives the topology for each grid. The common support lets the Gaussian and mesh errors be examined without changing the included pixels [\hyperlink{evE9}{E9}].

\begin{center}\begin{minipage}{\linewidth}
\small\begin{tabular*}{\linewidth}{@{\extracolsep{\fill}}llrrrr@{}}
\toprule
Output & Arm & Coverage & Signed median (cm) & Abs.\ median (cm) & Abs.\ P90 (cm)\\
\midrule
\inputrows{evidence/matte_grid_depth_rows.tex}
\bottomrule
\end{tabular*}\par
\captionof{table}[Matte depth discrepancy on shared support across grids]{Sensor-depth statistics on joint support across both Gaussian predictions and all six matte mesh grids. Coverage is measured against all stable sensor pixels.}\label{tab:grid-depth}
\end{minipage}\end{center}

\begin{center}\begin{minipage}{\linewidth}
\small\begin{tabular*}{\linewidth}{@{\extracolsep{\fill}}rlrrrr@{}}
\toprule
Grid & Arm & Triangles & Components & Boundary edges & Largest area\\
\midrule
\inputrows{evidence/matte_grid_topology_rows.tex}
\bottomrule
\end{tabular*}\par
\captionof{table}[Matte raw mesh topology across grid resolutions]{Raw mesh topology as grid resolution changes. Largest area is the fraction contributed by the largest edge-connected component by area.}\label{tab:grid-topology}
\end{minipage}\end{center}
\reportfigure[Matte depth discrepancy across TSDF grid resolutions]{0.99\linewidth}{figures/matte_grid_discrepancy.png}
  {Held-out matte depth discrepancy in frame rgb\_0006487.png. Blue means the prediction lies nearer to the camera than the aligned sensor depth. Light gray masks pixels outside the shared stable support. The displayed color range clips at $\pm20$ cm.}{fig:grid-depth}
The joint mask retains 856,383 of the 918,331 stable sensor pixels (93.25\%). On that mask, A05-Abs has 0.70 cm higher median absolute discrepancy before fusion. Its mesh has about 1.05--1.09 cm higher median absolute discrepancy at each tested grid. Changing grid resolution from 128 to 256 shifts each arm's median by less than 0.05 cm across the three settings. Thus some of the separation is present in the learned fields, and the measured difference survives the tested discretization change. Medians summarize the whole mask and do not decompose into additive per-pixel sources of error.

Grid spacing does matter for surface structure. In A05-Abs, coverage rises from 96.27\% to 98.65\% while components rise from 695 to 2,284 and boundary edges from 9,061 to 36,040. The largest connected region grows from 66.0\% to 92.1\% of surface area even as more small pieces appear. B00 shows rising coverage, components and boundaries too, with a more stable largest-area fraction. Every output retains open boundaries. Figure~\ref{fig:grid-depth} shows the spatial persistence of the negative depth discrepancy in an example frame. These results indicate substantial discretization dependence in topology alongside a more persistent depth offset. The common capture and pose pipeline limits physical attribution [\hyperlink{evE9}{E9}].


\subsection{Editing trials across scenes}
For each of the ten raw meshes, Blender selected a connected patch within four voxels of a corresponding ray hit. The same task displaced it toward the camera with a peak offset of half a voxel. All ten saved files reopened and retained displaced vertices. Both public Room meshes moved 108 vertices at the corresponding coffee-table patch. On matte, B00 moved 97 and A05-Abs moved 84. A selected patch can remain editable on a fragmented mesh. The test establishes local edit persistence at those locations [\hyperlink{evE4}{E4}].\par
\subsection{Public Room from raw mesh to design edit}
A separate Room experiment took the fixed A05-Abs raw PLY, removed small connected fragments satisfying both a face-count and an area threshold, then applied quadric decimation. Raw, filtered and cleaned meshes contained 119,987, 119,454 and 77,624 triangles respectively. The final mesh retained 748 components and 13,386 boundary edges, down from 1,162 and 16,807. Filtering and simplification reduced complexity while leaving open surfaces visible [\hyperlink{evE6}{E6}].\par
\reportfigure[Room silhouettes and mesh depth before and after cleanup]{0.86\linewidth}{figures/room_cleanup.png}
  {Public Room silhouette comparison at matched held-out camera crops. Each row shows one subject. Columns compare the photograph, raw and cleaned vertex colors, and raw and cleaned mesh depth. Inspect the broad table edge, the potted leaves and the thin footstool supports.}{figure12}
In the coffee-table crop of Figure~\ref{figure12}, cleanup changed zero hit pixels out of 96,292. The foliage crop lost 374 and gained 17. Thin footstool supports were already incomplete before simplification. A scripted Blender operation then lifted a connected 14-vertex patch on the front-left tabletop. Figure~\ref{figure13} shows the edited patch. The reopened shape key retained a peak displacement of 0.040 normalized scene units in +Z at unit object scale. Attribution of physical dimensions would require a metric reference for public Room [\hyperlink{evE6}{E6}].\par
\reportfigure[Tabletop patch before and after the Blender edit]{0.67\linewidth}{figures/tabletop_edit.jpg}
  {Before and after the tabletop edit, rendered in Blender from the same camera. The highlighted patch in the reopened file identifies the edited vertices. The logged vertex displacement confirms the geometric change independently of the highlight.}{figure13}
\clearpage
\subsection{Threats to validity}
\subsubsection{Pairing and attribution}
The input audit matches starting tensors, ordered fitting cameras and decoded targets. GPU accumulation can still vary with execution order, and the baseline's random split depends on its random draws. The Room factorial estimates the joint effects of gradient aggregation and threshold under one all-split policy and seed. The ranking study removes one cue at a time, with a binding cap providing competition. The archived A04 factorial estimates operator and small-parent effects under signed gradients. Interactions among these studies' factors remain unresolved.

\subsubsection{Evaluation exposure, sample size and performance}
The four public indoor validation sets were excluded from Gaussian fitting and inspected during development. Garden extends evaluation to an outdoor scene outside method selection. It shares the Mip-NeRF 360 dataset family, uses factor-four images, and has a single paired seed. Supplied SfM can use all benchmark photographs, including held-out RGB views. The three Room seeds measure sensitivity to initialization and sampling trajectories on the same development scene. The controlled Room factorial, tight-cap ranking ablation and matched-time comparison still have one seed per arm.

PSNR reduces RGB error to a scalar and can favor smooth estimates. In Bonsai the mean gain depended on one view, while the median gain was small and 17 of 37 targets regressed. The primary mean retains every target. The leave-one-view-out check and alternate metrics show what that mean obscures. Early learning curves include changing harmonic degree and representation size. Timings depend on hardware, loader settings and concurrent activity. The archived Bonsai pair had unequal worker counts. The controlled Room time comparison matched worker settings and excluded validation. Render-call throughput excludes display and interaction.

\subsubsection{Surface evidence and editor validity}
The three nearby matte evaluation views share calibration and pose estimates with the acquisition pipeline. Dense disagreement maps still describe three correlated observations, with possible motion, sensor or alignment bias. The stability filter removes much depth-edge evidence. Common support controls sample inclusion within the mask, and coverage reports what falls outside each output. The grid study tests discretization with fixed crop, depth input and truncation. Similar tests with varied render resolution or depth statistics remain open.

Open boundaries and vertex-manifoldness failures constrain downstream modeling. Component counts mix separated objects with fragments. The report uses \emph{editable} for the observed connected-patch operation and save/reopen persistence. Self-intersections, dimensional tolerances and complex modeling operations require separate verification.

\clearpage
\section{Conclusion and Future Work}
\subsection{Representation efficiency}
Across the four indoor development scenes, A05-Abs used 45.57--54.56\% fewer Gaussians and gained 0.210--0.739 dB at equal updates. Removing Bonsai's largest gain changed its quantized mean difference from $+0.337$ to $-0.023$ dB. In Room the selected 452,409-Gaussian model scored 30.757 dB. A lower threshold bought higher PSNR at a larger count, and a binding 400,000 cap gave a smaller model at 30.706 dB. The selected variant also led PSNR at the three common time probes [\hyperlink{evE4}{E4}, \hyperlink{evE8}{E8}, \hyperlink{evE11}{E11}].

The Room gain repeated across three seeds. In Garden, the frozen policy gained 0.365 dB with 75.6\% fewer primitives and a binding 700,000 ceiling. This extends the measured trade-off to one outdoor scene in the same benchmark family [\hyperlink{evE12}{E12}].

\subsection{Allocation and implementation}
The selected Room run admitted every eligible parent. The gradient and split studies support the choice of eligibility rule and split operator, although their interaction under absolute gradients remains open. At the tighter 400,000 cap, removing either ranking cue changed which parents could split and produced mixed RGB results. The study therefore retains the full policy as its reference [\hyperlink{evE3}{E3}, \hyperlink{evE8}{E8}].

\subsection{Surface reconstruction and editing}
All ten primary benchmark arms yielded polygon meshes. On Room, the depth-curvature penalty reduced the candidate's curvature proxy by 3.06\% but increased the number of mesh components from 1,162 to 1,216. The final candidate therefore omits the penalty [\hyperlink{evE5}{E5}]. Room and Bonsai had lower raw component and boundary counts with the candidate, while Kitchen and the museum matte segment showed the reverse. The matte mesh covered slightly more stable depth, yet its median absolute sensor discrepancy rose from 8.38 to 9.47 cm on the original paired mesh support. The median discrepancy was already larger before fusion and persisted across all three grid resolutions. Physical surface accuracy remains unverified [\hyperlink{evE4}{E4}, \hyperlink{evE9}{E9}].

All ten matched raw-mesh edits persisted after reopening. The separate cleaned public Room demonstration and the fresh replay both preserved the tabletop lift in Blender, at $0.040$ normalized units. Open boundaries remain in the edited asset. These tests establish local edit persistence [\hyperlink{evE6}{E6}, \hyperlink{evE10}{E10}].

\subsection{Follow-up experiments}
Run the frozen policy on a second dataset family and vary the split operator and ranking jointly under absolute gradients. Repeated measurements on another GPU would quantify hardware-dependent training time.

Evaluate completeness and dimensional error against surveyed geometry with known scale, including depth edges filtered from the matte comparison. Vary render resolution, fusion truncation and crop independently at fixed checkpoints. Verify additional editing operations with explicit dimensional tolerances.

\clearpage
\phantomsection
\section*{Appendix A. Reproduce the Public Room Pipeline}
\addcontentsline{toc}{section}{Appendix A. Reproduce the Public Room Pipeline}
Extract \path{final_report_supplement.zip} and open PowerShell at its root. Configure Python~3.12, PyTorch, CUDA~12.8 and MSVC in an x64 developer shell by following \path{report/final_report/REPRODUCE.md}. The public data come from the official Mip-NeRF 360 archive. The supplement supplies the project source, the fitting order and the verified Room edit [\hyperlink{evE10}{E10}].

The upstream archive was fetched again into an empty preparation directory. The fresh preparation matched all 311 existing factor-two photo files, the three original COLMAP binaries and the five files in the prepared sparse model byte for byte. The ordered fitting indices are shipped with the report. The benchmark dataset audit supplies the fitting and held-out identities [\hyperlink{evE1}{E1}, \hyperlink{evE10}{E10}].

\subsection*{A.1. Run the public Room pipeline}
Use the commands below in order. Adjust the Blender path if needed and choose a fresh results directory. The starter script stores the exact native CLI flags for each stage. The equivalent one-command run uses \texttt{-Stage All}.

{\small
\begin{verbatim}
$run = '.\report\final_report\reproduce_room.ps1'
$opts = @{
    Python = '.\.venv\Scripts\python.exe'
    Blender = 'C:\Program Files\Blender Foundation\Blender 4.5\blender.exe'
    Work = 'results/report_room_replay'
}
& $run @opts -Stage Prepare
& $run @opts -Stage Train
& $run @opts -Stage Extract
& $run @opts -Stage Clean
& $run @opts -Stage Edit
& $run @opts -Stage Verify
\end{verbatim}
}
Table~\ref{tab:replay-artifacts} lists the expected stage outputs. Training uses 7,000 updates, seed 42, one packed image per batch and the frozen fitting order. Extraction uses the post-update checkpoint, 64 fitting cameras, quarter-resolution renders and a uniform TSDF grid up to 192 cells per axis. Verify reopens the saved project in a separate Blender process with autoexecution disabled and measures the tabletop patch in the evaluated mesh.

\begin{center}\begin{minipage}{\linewidth}
\small\begin{tabularx}{\linewidth}{@{}p{.17\linewidth}YY@{}}
\toprule
Stage & Files & Expected check\\
\midrule
Preparation & \texttt{dataset\_manifest.json}, \texttt{images/}, \texttt{sparse/0/} & 311 input images, 39 validation views, 112,627 sparse points and matching intrinsics\\
Training & \texttt{experiment.json}, \texttt{final\_post\_update.pt}, trainer statistics and configuration & 7,000 updates, finite tensors, matched initialization and complete camera order\\
Extraction & \texttt{raw\_mesh.ply}, \texttt{raw\_mesh.obj}, \texttt{report.json} & Nonempty triangles and matching camera, checkpoint and mesh hashes\\
Cleanup & \texttt{room\_cleaned.ply}, \texttt{cleanup.json}, camera crops & Verified raw input and reloaded derivative topology\\
Editing & \texttt{room\_tabletop\_edited.blend}, \texttt{design\_reopen.json} & Verified saved hash and evaluated displacement in a separate process\\
\bottomrule
\end{tabularx}\par
\captionof{table}[Expected outputs and checks for a complete Room run]{Files that identify a complete public Room run.}\label{tab:replay-artifacts}
\end{minipage}\end{center}

\clearpage
The fresh Room rehearsal used the verified 452,409-Gaussian checkpoint from the matched replay. The uniform TSDF produced 119,061 triangles. Cleanup reduced that to 76,971 triangles, with 779 components and 13,559 open boundary edges. Blender reopened the separate project and confirmed a $+0.040$ peak lift in normalized units on a connected 12-vertex patch. The archived demonstration used 14 vertices. The repeated operation passed on the new surface with a different triangulation. The run record links all input and output digests [\hyperlink{evE10}{E10}].


\subsection*{A.2. Replay the controlled studies}
The original controlled Room runner uses historical baseline target panels for its replay audit. The supplement contains the frozen protocol and verified results. The public standalone procedure above starts from prepared photographs and poses. The additional paired evaluation can also run from fresh benchmark inputs. Its commands and required scenes are listed in the guide.

The matte sweep requires the original paired Week 12 checkpoints, prepared local data, extraction records and the three reserved depth messages. It verifies all relevant hashes. It uses the same support mask across the six meshes and two Gaussian predictions. Its JSON record retains all per-view depth summaries, topology and coverage. The local captures remain under their documented project paths. Section~\ref{sec:grid-results} identifies the measured comparison and its limits.

\subsection*{A.3. Inspect common failures}
An initialization hash mismatch points to a different input model, parser or runtime. A camera-order mismatch can arise despite matching seeds if image IDs, loading or random streams change. Start by checking the exact dataset identities and replay file. A checkpoint with nonfinite parameters is invalid even if a subset of rendered images appears normal. Empty TSDF output may indicate a camera-convention, volume-bounds or valid-depth problem. Inspect those inputs before loosening the thresholds.

If the cleaned Room topology or imported bounds differ from its recorded source, inspect the PLY digest and vertex import scale. If Blender cannot recover a connected tabletop patch, preserve that failed run and inspect the camera crop and reopened displacement record. A screenshot of the selected vertices cannot substitute for the measured displacement. Benchmark scores and mesh counts can differ across runs despite a matching input protocol, so record observed values and identify any changed versions or conditions.


\clearpage
\phantomsection
\section*{Evidence Register}
\addcontentsline{toc}{section}{Evidence Register}
The IDs below are stable destinations within this PDF. Unzip \texttt{final\_report\_supplement.zip} and open \texttt{EVIDENCE\_INDEX.md} at its root for the exact files keyed by these IDs. \texttt{MANIFEST.sha256} provides a digest for each bundled file. The supplement includes the public Room code and a saved edit, the reported numeric evidence and the reproducible plots. Public photographs are retrieved from the benchmark. The original museum recording remains with the project.

\hypertarget{evE1}{\textbf{[E1]}} Dataset dimensions, camera counts and fitting/held-out identities for the public indoor benchmark.

\hypertarget{evE2}{\textbf{[E2]}} Local RGB-D fitting, withheld-view localization and the screening of the complete museum route.

\hypertarget{evE3}{\textbf{[E3]}} Strategy implementation, earlier multi-seed campaign, split operator factorial and absolute-gradient studies.

\hypertarget{evE4}{\textbf{[E4]}} Matched Room, Counter, Kitchen, Bonsai and matte pair records, with trainer RGB metrics, fitting sequence, checkpoint digests and surface extraction summaries. The original 7,000-update comparison uses evaluation index 6999.

\hypertarget{evE5}{\textbf{[E5]}} Surface-curvature ablation, replay correction and known-plane control.

\hypertarget{evE6}{\textbf{[E6]}} Original Room cleanup and Blender reopen measurements, supporting code and the newer bundled editable asset.

\hypertarget{evE7}{\textbf{[E7]}} Progress report and the secondary ordinary-video acquisition workflow.

\hypertarget{evE8}{\textbf{[E8]}} Controlled Room replays with verified endpoint and timed outcomes, exact fitting order, cap sweep and tight-cap ranking protocol.

\hypertarget{evE9}{\textbf{[E9]}} Fixed-checkpoint matte grid study and shared-support depth measurements.

\hypertarget{evE10}{\textbf{[E10]}} Public Room reproduction commands, fresh input audit, raw geometry, cleanup and independently reopened Blender asset.

\hypertarget{evE11}{\textbf{[E11]}} Archived mesh topology and component areas, and paired per-view RGB evidence.

\hypertarget{evE12}{\textbf{[E12]}} Frozen Garden and Room seed protocol, verified paired results and Garden per-view RGB checks.

\clearpage\phantomsection\addcontentsline{toc}{section}{References}
\begin{thebibliography}{15}\small\setlength{\itemsep}{5pt}
\bibitem{ref1} J. L. Schönberger and J.-M. Frahm, \reftitle{https://openaccess.thecvf.com/content_cvpr_2016/html/Schonberger_Structure-From-Motion_Revisited_CVPR_2016_paper.html}{Structure-from-motion revisited}, in \emph{Proc. IEEE Conf. Computer Vision and Pattern Recognition (CVPR)}, 2016, pp. 4104--4113.
\bibitem{ref2} J. T. Barron, B. Mildenhall, D. Verbin, P. P. Srinivasan, and P. Hedman, \reftitle{https://openaccess.thecvf.com/content/CVPR2022/html/Barron_Mip-NeRF_360_Unbounded_Anti-Aliased_Neural_Radiance_Fields_CVPR_2022_paper.html}{Mip-NeRF 360: Unbounded anti-aliased neural radiance fields}, in \emph{Proc. IEEE/CVF Conf. Computer Vision and Pattern Recognition (CVPR)}, 2022, pp. 5470--5479.
\bibitem{ref3} B. Mildenhall, P. P. Srinivasan, M. Tancik, J. T. Barron, R. Ramamoorthi, and R. Ng, \reftitle{https://doi.org/10.1007/978-3-030-58452-8_24}{NeRF: Representing scenes as neural radiance fields for view synthesis}, in \emph{Proc. European Conf. Computer Vision (ECCV)}, 2020, pp. 405--421.
\bibitem{ref4} B. Kerbl, G. Kopanas, T. Leimkühler, and G. Drettakis, \reftitle{https://doi.org/10.1145/3592433}{3D Gaussian splatting for real-time radiance field rendering}, \emph{ACM Trans. Graphics}, vol. 42, no. 4, Art. no. 139, 2023.
\bibitem{ref5} V. Ye, R. Li, J. Kerr, M. Turkulainen, B. Yi, Z. Pan, O. Seiskari, J. Ye, J. Hu, M. Tancik, and A. Kanazawa, \reftitle{https://www.jmlr.org/papers/v26/24-1476.html}{gsplat: An open-source library for Gaussian splatting}, \emph{Journal of Machine Learning Research}, vol. 26, no. 34, pp. 1--17, 2025.
\bibitem{ref6} Z. Ye, W. Li, S. Liu, P. Qiao, and Y. Dou, \reftitle{https://doi.org/10.1145/3664647.3681361}{AbsGS: Recovering fine details in 3D Gaussian splatting}, in \emph{Proc. ACM Int. Conf. Multimedia (MM)}, 2024, pp. 1053--1061.
\bibitem{ref7} X. Deng, C. Diao, M. Li, R. Yu, and D. Xu, \reftitle{https://openaccess.thecvf.com/content/CVPR2026F/html/Deng_Improving_Densification_in_3D_Gaussian_Splatting_for_High-Fidelity_Rendering_CVPRF_2026_paper.html}{Improving densification in 3D Gaussian splatting for high-fidelity rendering}, in \emph{Proc. IEEE/CVF Conf. Computer Vision and Pattern Recognition Findings (CVPR Findings)}, 2026, pp. 223--232.
\bibitem{ref8} A. Guédon and V. Lepetit, \reftitle{https://openaccess.thecvf.com/content/CVPR2024/html/Guedon_SuGaR_Surface-Aligned_Gaussian_Splatting_for_Efficient_3D_Mesh_Reconstruction_and_CVPR_2024_paper.html}{SuGaR: Surface-aligned Gaussian splatting for efficient 3D mesh reconstruction and high-quality mesh rendering}, in \emph{Proc. IEEE/CVF Conf. Computer Vision and Pattern Recognition (CVPR)}, 2024, pp. 5354--5363.
\bibitem{ref9} B. Curless and M. Levoy, \reftitle{https://www.graphics.stanford.edu/papers/volrange/}{A volumetric method for building complex models from range images}, in \emph{Proc. ACM SIGGRAPH}, 1996, pp. 303--312.
\bibitem{ref10} Z. Wang, A. C. Bovik, H. R. Sheikh, and E. P. Simoncelli, \reftitle{https://doi.org/10.1109/TIP.2003.819861}{Image quality assessment: From error visibility to structural similarity}, \emph{IEEE Trans. Image Processing}, vol. 13, no. 4, pp. 600--612, 2004.
\bibitem{ref11} R. Zhang, P. Isola, A. A. Efros, E. Shechtman, and O. Wang, \reftitle{https://openaccess.thecvf.com/content_cvpr_2018/html/Zhang_The_Unreasonable_Effectiveness_CVPR_2018_paper.html}{The unreasonable effectiveness of deep features as a perceptual metric}, in \emph{Proc. IEEE/CVF Conf. Computer Vision and Pattern Recognition (CVPR)}, 2018, pp. 586--595.
\bibitem{ref12} S. S. Mallick, R. Goel, B. Kerbl, M. Steinberger, F. V. Carrasco, and F. De la Torre, \reftitle{https://doi.org/10.1145/3680528.3687694}{Taming 3DGS: High-quality radiance fields with limited resources}, in \emph{Proc. ACM SIGGRAPH Asia}, 2024, Art. no. 2.
\bibitem{ref13} B. Huang, Z. Yu, A. Chen, A. Geiger, and S. Gao, \reftitle{https://doi.org/10.1145/3641519.3657428}{2D Gaussian splatting for geometrically accurate radiance fields}, in \emph{Proc. ACM SIGGRAPH}, 2024, Art. no. 32.
\bibitem{ref14} W. E. Lorensen and H. E. Cline, \reftitle{https://doi.org/10.1145/37402.37422}{Marching cubes: A high resolution 3D surface construction algorithm}, \emph{ACM SIGGRAPH Computer Graphics}, vol. 21, no. 4, pp. 163--169, 1987.
\bibitem{ref15} M. Kazhdan, M. Bolitho, and H. Hoppe, \reftitle{https://diglib.eg.org/items/245d6ba1-0209-4365-a4a4-16ce46421620}{Poisson surface reconstruction}, in \emph{Proc. Eurographics Symposium on Geometry Processing}, 2006, pp. 61--70.
\end{thebibliography}

\label{report-end}
\end{document}
